\documentclass[11pt]{article}

    \usepackage[breakable]{tcolorbox}
    \usepackage{parskip} % Stop auto-indenting (to mimic markdown behaviour)
    

    % Basic figure setup, for now with no caption control since it's done
    % automatically by Pandoc (which extracts ![](path) syntax from Markdown).
    \usepackage{graphicx}
    % Keep aspect ratio if custom image width or height is specified
    \setkeys{Gin}{keepaspectratio}
    % Maintain compatibility with old templates. Remove in nbconvert 6.0
    \let\Oldincludegraphics\includegraphics
    % Ensure that by default, figures have no caption (until we provide a
    % proper Figure object with a Caption API and a way to capture that
    % in the conversion process - todo).
    \usepackage{caption}
    \DeclareCaptionFormat{nocaption}{}
    \captionsetup{format=nocaption,aboveskip=0pt,belowskip=0pt}

    \usepackage{float}
    \floatplacement{figure}{H} % forces figures to be placed at the correct location
    \usepackage{xcolor} % Allow colors to be defined
    \usepackage{enumerate} % Needed for markdown enumerations to work
    \usepackage{geometry} % Used to adjust the document margins
    \usepackage{amsmath} % Equations
    \usepackage{amssymb} % Equations
    \usepackage{textcomp} % defines textquotesingle
    % Hack from http://tex.stackexchange.com/a/47451/13684:
    \AtBeginDocument{%
        \def\PYZsq{\textquotesingle}% Upright quotes in Pygmentized code
    }
    \usepackage{upquote} % Upright quotes for verbatim code
    \usepackage{eurosym} % defines \euro

    \usepackage{iftex}
    \ifPDFTeX
        \usepackage[T1]{fontenc}
        \IfFileExists{alphabeta.sty}{
              \usepackage{alphabeta}
          }{
              \usepackage[mathletters]{ucs}
              \usepackage[utf8x]{inputenc}
          }
    \else
        \usepackage{fontspec}
        \usepackage{unicode-math}
    \fi

    \usepackage{fancyvrb} % verbatim replacement that allows latex
    \usepackage{grffile} % extends the file name processing of package graphics
                         % to support a larger range
    \makeatletter % fix for old versions of grffile with XeLaTeX
    \@ifpackagelater{grffile}{2019/11/01}
    {
      % Do nothing on new versions
    }
    {
      \def\Gread@@xetex#1{%
        \IfFileExists{"\Gin@base".bb}%
        {\Gread@eps{\Gin@base.bb}}%
        {\Gread@@xetex@aux#1}%
      }
    }
    \makeatother
    \usepackage[Export]{adjustbox} % Used to constrain images to a maximum size
    \adjustboxset{max size={0.9\linewidth}{0.9\paperheight}}

    % The hyperref package gives us a pdf with properly built
    % internal navigation ('pdf bookmarks' for the table of contents,
    % internal cross-reference links, web links for URLs, etc.)
    \usepackage{hyperref}
    % The default LaTeX title has an obnoxious amount of whitespace. By default,
    % titling removes some of it. It also provides customization options.
    \usepackage{titling}
    \usepackage{longtable} % longtable support required by pandoc >1.10
    \usepackage{booktabs}  % table support for pandoc > 1.12.2
    \usepackage{array}     % table support for pandoc >= 2.11.3
    \usepackage{calc}      % table minipage width calculation for pandoc >= 2.11.1
    \usepackage[inline]{enumitem} % IRkernel/repr support (it uses the enumerate* environment)
    \usepackage[normalem]{ulem} % ulem is needed to support strikethroughs (\sout)
                                % normalem makes italics be italics, not underlines
    \usepackage{soul}      % strikethrough (\st) support for pandoc >= 3.0.0
    \usepackage{mathrsfs}
    

    
    % Colors for the hyperref package
    \definecolor{urlcolor}{rgb}{0,.145,.698}
    \definecolor{linkcolor}{rgb}{.71,0.21,0.01}
    \definecolor{citecolor}{rgb}{.12,.54,.11}

    % ANSI colors
    \definecolor{ansi-black}{HTML}{3E424D}
    \definecolor{ansi-black-intense}{HTML}{282C36}
    \definecolor{ansi-red}{HTML}{E75C58}
    \definecolor{ansi-red-intense}{HTML}{B22B31}
    \definecolor{ansi-green}{HTML}{00A250}
    \definecolor{ansi-green-intense}{HTML}{007427}
    \definecolor{ansi-yellow}{HTML}{DDB62B}
    \definecolor{ansi-yellow-intense}{HTML}{B27D12}
    \definecolor{ansi-blue}{HTML}{208FFB}
    \definecolor{ansi-blue-intense}{HTML}{0065CA}
    \definecolor{ansi-magenta}{HTML}{D160C4}
    \definecolor{ansi-magenta-intense}{HTML}{A03196}
    \definecolor{ansi-cyan}{HTML}{60C6C8}
    \definecolor{ansi-cyan-intense}{HTML}{258F8F}
    \definecolor{ansi-white}{HTML}{C5C1B4}
    \definecolor{ansi-white-intense}{HTML}{A1A6B2}
    \definecolor{ansi-default-inverse-fg}{HTML}{FFFFFF}
    \definecolor{ansi-default-inverse-bg}{HTML}{000000}

    % common color for the border for error outputs.
    \definecolor{outerrorbackground}{HTML}{FFDFDF}

    % commands and environments needed by pandoc snippets
    % extracted from the output of `pandoc -s`
    \providecommand{\tightlist}{%
      \setlength{\itemsep}{0pt}\setlength{\parskip}{0pt}}
    \DefineVerbatimEnvironment{Highlighting}{Verbatim}{commandchars=\\\{\}}
    % Add ',fontsize=\small' for more characters per line
    \newenvironment{Shaded}{}{}
    \newcommand{\KeywordTok}[1]{\textcolor[rgb]{0.00,0.44,0.13}{\textbf{{#1}}}}
    \newcommand{\DataTypeTok}[1]{\textcolor[rgb]{0.56,0.13,0.00}{{#1}}}
    \newcommand{\DecValTok}[1]{\textcolor[rgb]{0.25,0.63,0.44}{{#1}}}
    \newcommand{\BaseNTok}[1]{\textcolor[rgb]{0.25,0.63,0.44}{{#1}}}
    \newcommand{\FloatTok}[1]{\textcolor[rgb]{0.25,0.63,0.44}{{#1}}}
    \newcommand{\CharTok}[1]{\textcolor[rgb]{0.25,0.44,0.63}{{#1}}}
    \newcommand{\StringTok}[1]{\textcolor[rgb]{0.25,0.44,0.63}{{#1}}}
    \newcommand{\CommentTok}[1]{\textcolor[rgb]{0.38,0.63,0.69}{\textit{{#1}}}}
    \newcommand{\OtherTok}[1]{\textcolor[rgb]{0.00,0.44,0.13}{{#1}}}
    \newcommand{\AlertTok}[1]{\textcolor[rgb]{1.00,0.00,0.00}{\textbf{{#1}}}}
    \newcommand{\FunctionTok}[1]{\textcolor[rgb]{0.02,0.16,0.49}{{#1}}}
    \newcommand{\RegionMarkerTok}[1]{{#1}}
    \newcommand{\ErrorTok}[1]{\textcolor[rgb]{1.00,0.00,0.00}{\textbf{{#1}}}}
    \newcommand{\NormalTok}[1]{{#1}}

    % Additional commands for more recent versions of Pandoc
    \newcommand{\ConstantTok}[1]{\textcolor[rgb]{0.53,0.00,0.00}{{#1}}}
    \newcommand{\SpecialCharTok}[1]{\textcolor[rgb]{0.25,0.44,0.63}{{#1}}}
    \newcommand{\VerbatimStringTok}[1]{\textcolor[rgb]{0.25,0.44,0.63}{{#1}}}
    \newcommand{\SpecialStringTok}[1]{\textcolor[rgb]{0.73,0.40,0.53}{{#1}}}
    \newcommand{\ImportTok}[1]{{#1}}
    \newcommand{\DocumentationTok}[1]{\textcolor[rgb]{0.73,0.13,0.13}{\textit{{#1}}}}
    \newcommand{\AnnotationTok}[1]{\textcolor[rgb]{0.38,0.63,0.69}{\textbf{\textit{{#1}}}}}
    \newcommand{\CommentVarTok}[1]{\textcolor[rgb]{0.38,0.63,0.69}{\textbf{\textit{{#1}}}}}
    \newcommand{\VariableTok}[1]{\textcolor[rgb]{0.10,0.09,0.49}{{#1}}}
    \newcommand{\ControlFlowTok}[1]{\textcolor[rgb]{0.00,0.44,0.13}{\textbf{{#1}}}}
    \newcommand{\OperatorTok}[1]{\textcolor[rgb]{0.40,0.40,0.40}{{#1}}}
    \newcommand{\BuiltInTok}[1]{{#1}}
    \newcommand{\ExtensionTok}[1]{{#1}}
    \newcommand{\PreprocessorTok}[1]{\textcolor[rgb]{0.74,0.48,0.00}{{#1}}}
    \newcommand{\AttributeTok}[1]{\textcolor[rgb]{0.49,0.56,0.16}{{#1}}}
    \newcommand{\InformationTok}[1]{\textcolor[rgb]{0.38,0.63,0.69}{\textbf{\textit{{#1}}}}}
    \newcommand{\WarningTok}[1]{\textcolor[rgb]{0.38,0.63,0.69}{\textbf{\textit{{#1}}}}}
    \makeatletter
    \newsavebox\pandoc@box
    \newcommand*\pandocbounded[1]{%
      \sbox\pandoc@box{#1}%
      % scaling factors for width and height
      \Gscale@div\@tempa\textheight{\dimexpr\ht\pandoc@box+\dp\pandoc@box\relax}%
      \Gscale@div\@tempb\linewidth{\wd\pandoc@box}%
      % select the smaller of both
      \ifdim\@tempb\p@<\@tempa\p@
        \let\@tempa\@tempb
      \fi
      % scaling accordingly (\@tempa < 1)
      \ifdim\@tempa\p@<\p@
        \scalebox{\@tempa}{\usebox\pandoc@box}%
      % scaling not needed, use as it is
      \else
        \usebox{\pandoc@box}%
      \fi
    }
    \makeatother

    % Define a nice break command that doesn't care if a line doesn't already
    % exist.
    \def\br{\hspace*{\fill} \\* }
    % Math Jax compatibility definitions
    \def\gt{>}
    \def\lt{<}
    \let\Oldtex\TeX
    \let\Oldlatex\LaTeX
    \renewcommand{\TeX}{\textrm{\Oldtex}}
    \renewcommand{\LaTeX}{\textrm{\Oldlatex}}
    % Document parameters
    % Document title
    \title{Ch03-linreg-lab}
    
    
    
    
    
    
    
% Pygments definitions
\makeatletter
\def\PY@reset{\let\PY@it=\relax \let\PY@bf=\relax%
    \let\PY@ul=\relax \let\PY@tc=\relax%
    \let\PY@bc=\relax \let\PY@ff=\relax}
\def\PY@tok#1{\csname PY@tok@#1\endcsname}
\def\PY@toks#1+{\ifx\relax#1\empty\else%
    \PY@tok{#1}\expandafter\PY@toks\fi}
\def\PY@do#1{\PY@bc{\PY@tc{\PY@ul{%
    \PY@it{\PY@bf{\PY@ff{#1}}}}}}}
\def\PY#1#2{\PY@reset\PY@toks#1+\relax+\PY@do{#2}}

\@namedef{PY@tok@w}{\def\PY@tc##1{\textcolor[rgb]{0.73,0.73,0.73}{##1}}}
\@namedef{PY@tok@c}{\let\PY@it=\textit\def\PY@tc##1{\textcolor[rgb]{0.24,0.48,0.48}{##1}}}
\@namedef{PY@tok@cp}{\def\PY@tc##1{\textcolor[rgb]{0.61,0.40,0.00}{##1}}}
\@namedef{PY@tok@k}{\let\PY@bf=\textbf\def\PY@tc##1{\textcolor[rgb]{0.00,0.50,0.00}{##1}}}
\@namedef{PY@tok@kp}{\def\PY@tc##1{\textcolor[rgb]{0.00,0.50,0.00}{##1}}}
\@namedef{PY@tok@kt}{\def\PY@tc##1{\textcolor[rgb]{0.69,0.00,0.25}{##1}}}
\@namedef{PY@tok@o}{\def\PY@tc##1{\textcolor[rgb]{0.40,0.40,0.40}{##1}}}
\@namedef{PY@tok@ow}{\let\PY@bf=\textbf\def\PY@tc##1{\textcolor[rgb]{0.67,0.13,1.00}{##1}}}
\@namedef{PY@tok@nb}{\def\PY@tc##1{\textcolor[rgb]{0.00,0.50,0.00}{##1}}}
\@namedef{PY@tok@nf}{\def\PY@tc##1{\textcolor[rgb]{0.00,0.00,1.00}{##1}}}
\@namedef{PY@tok@nc}{\let\PY@bf=\textbf\def\PY@tc##1{\textcolor[rgb]{0.00,0.00,1.00}{##1}}}
\@namedef{PY@tok@nn}{\let\PY@bf=\textbf\def\PY@tc##1{\textcolor[rgb]{0.00,0.00,1.00}{##1}}}
\@namedef{PY@tok@ne}{\let\PY@bf=\textbf\def\PY@tc##1{\textcolor[rgb]{0.80,0.25,0.22}{##1}}}
\@namedef{PY@tok@nv}{\def\PY@tc##1{\textcolor[rgb]{0.10,0.09,0.49}{##1}}}
\@namedef{PY@tok@no}{\def\PY@tc##1{\textcolor[rgb]{0.53,0.00,0.00}{##1}}}
\@namedef{PY@tok@nl}{\def\PY@tc##1{\textcolor[rgb]{0.46,0.46,0.00}{##1}}}
\@namedef{PY@tok@ni}{\let\PY@bf=\textbf\def\PY@tc##1{\textcolor[rgb]{0.44,0.44,0.44}{##1}}}
\@namedef{PY@tok@na}{\def\PY@tc##1{\textcolor[rgb]{0.41,0.47,0.13}{##1}}}
\@namedef{PY@tok@nt}{\let\PY@bf=\textbf\def\PY@tc##1{\textcolor[rgb]{0.00,0.50,0.00}{##1}}}
\@namedef{PY@tok@nd}{\def\PY@tc##1{\textcolor[rgb]{0.67,0.13,1.00}{##1}}}
\@namedef{PY@tok@s}{\def\PY@tc##1{\textcolor[rgb]{0.73,0.13,0.13}{##1}}}
\@namedef{PY@tok@sd}{\let\PY@it=\textit\def\PY@tc##1{\textcolor[rgb]{0.73,0.13,0.13}{##1}}}
\@namedef{PY@tok@si}{\let\PY@bf=\textbf\def\PY@tc##1{\textcolor[rgb]{0.64,0.35,0.47}{##1}}}
\@namedef{PY@tok@se}{\let\PY@bf=\textbf\def\PY@tc##1{\textcolor[rgb]{0.67,0.36,0.12}{##1}}}
\@namedef{PY@tok@sr}{\def\PY@tc##1{\textcolor[rgb]{0.64,0.35,0.47}{##1}}}
\@namedef{PY@tok@ss}{\def\PY@tc##1{\textcolor[rgb]{0.10,0.09,0.49}{##1}}}
\@namedef{PY@tok@sx}{\def\PY@tc##1{\textcolor[rgb]{0.00,0.50,0.00}{##1}}}
\@namedef{PY@tok@m}{\def\PY@tc##1{\textcolor[rgb]{0.40,0.40,0.40}{##1}}}
\@namedef{PY@tok@gh}{\let\PY@bf=\textbf\def\PY@tc##1{\textcolor[rgb]{0.00,0.00,0.50}{##1}}}
\@namedef{PY@tok@gu}{\let\PY@bf=\textbf\def\PY@tc##1{\textcolor[rgb]{0.50,0.00,0.50}{##1}}}
\@namedef{PY@tok@gd}{\def\PY@tc##1{\textcolor[rgb]{0.63,0.00,0.00}{##1}}}
\@namedef{PY@tok@gi}{\def\PY@tc##1{\textcolor[rgb]{0.00,0.52,0.00}{##1}}}
\@namedef{PY@tok@gr}{\def\PY@tc##1{\textcolor[rgb]{0.89,0.00,0.00}{##1}}}
\@namedef{PY@tok@ge}{\let\PY@it=\textit}
\@namedef{PY@tok@gs}{\let\PY@bf=\textbf}
\@namedef{PY@tok@ges}{\let\PY@bf=\textbf\let\PY@it=\textit}
\@namedef{PY@tok@gp}{\let\PY@bf=\textbf\def\PY@tc##1{\textcolor[rgb]{0.00,0.00,0.50}{##1}}}
\@namedef{PY@tok@go}{\def\PY@tc##1{\textcolor[rgb]{0.44,0.44,0.44}{##1}}}
\@namedef{PY@tok@gt}{\def\PY@tc##1{\textcolor[rgb]{0.00,0.27,0.87}{##1}}}
\@namedef{PY@tok@err}{\def\PY@bc##1{{\setlength{\fboxsep}{\string -\fboxrule}\fcolorbox[rgb]{1.00,0.00,0.00}{1,1,1}{\strut ##1}}}}
\@namedef{PY@tok@kc}{\let\PY@bf=\textbf\def\PY@tc##1{\textcolor[rgb]{0.00,0.50,0.00}{##1}}}
\@namedef{PY@tok@kd}{\let\PY@bf=\textbf\def\PY@tc##1{\textcolor[rgb]{0.00,0.50,0.00}{##1}}}
\@namedef{PY@tok@kn}{\let\PY@bf=\textbf\def\PY@tc##1{\textcolor[rgb]{0.00,0.50,0.00}{##1}}}
\@namedef{PY@tok@kr}{\let\PY@bf=\textbf\def\PY@tc##1{\textcolor[rgb]{0.00,0.50,0.00}{##1}}}
\@namedef{PY@tok@bp}{\def\PY@tc##1{\textcolor[rgb]{0.00,0.50,0.00}{##1}}}
\@namedef{PY@tok@fm}{\def\PY@tc##1{\textcolor[rgb]{0.00,0.00,1.00}{##1}}}
\@namedef{PY@tok@vc}{\def\PY@tc##1{\textcolor[rgb]{0.10,0.09,0.49}{##1}}}
\@namedef{PY@tok@vg}{\def\PY@tc##1{\textcolor[rgb]{0.10,0.09,0.49}{##1}}}
\@namedef{PY@tok@vi}{\def\PY@tc##1{\textcolor[rgb]{0.10,0.09,0.49}{##1}}}
\@namedef{PY@tok@vm}{\def\PY@tc##1{\textcolor[rgb]{0.10,0.09,0.49}{##1}}}
\@namedef{PY@tok@sa}{\def\PY@tc##1{\textcolor[rgb]{0.73,0.13,0.13}{##1}}}
\@namedef{PY@tok@sb}{\def\PY@tc##1{\textcolor[rgb]{0.73,0.13,0.13}{##1}}}
\@namedef{PY@tok@sc}{\def\PY@tc##1{\textcolor[rgb]{0.73,0.13,0.13}{##1}}}
\@namedef{PY@tok@dl}{\def\PY@tc##1{\textcolor[rgb]{0.73,0.13,0.13}{##1}}}
\@namedef{PY@tok@s2}{\def\PY@tc##1{\textcolor[rgb]{0.73,0.13,0.13}{##1}}}
\@namedef{PY@tok@sh}{\def\PY@tc##1{\textcolor[rgb]{0.73,0.13,0.13}{##1}}}
\@namedef{PY@tok@s1}{\def\PY@tc##1{\textcolor[rgb]{0.73,0.13,0.13}{##1}}}
\@namedef{PY@tok@mb}{\def\PY@tc##1{\textcolor[rgb]{0.40,0.40,0.40}{##1}}}
\@namedef{PY@tok@mf}{\def\PY@tc##1{\textcolor[rgb]{0.40,0.40,0.40}{##1}}}
\@namedef{PY@tok@mh}{\def\PY@tc##1{\textcolor[rgb]{0.40,0.40,0.40}{##1}}}
\@namedef{PY@tok@mi}{\def\PY@tc##1{\textcolor[rgb]{0.40,0.40,0.40}{##1}}}
\@namedef{PY@tok@il}{\def\PY@tc##1{\textcolor[rgb]{0.40,0.40,0.40}{##1}}}
\@namedef{PY@tok@mo}{\def\PY@tc##1{\textcolor[rgb]{0.40,0.40,0.40}{##1}}}
\@namedef{PY@tok@ch}{\let\PY@it=\textit\def\PY@tc##1{\textcolor[rgb]{0.24,0.48,0.48}{##1}}}
\@namedef{PY@tok@cm}{\let\PY@it=\textit\def\PY@tc##1{\textcolor[rgb]{0.24,0.48,0.48}{##1}}}
\@namedef{PY@tok@cpf}{\let\PY@it=\textit\def\PY@tc##1{\textcolor[rgb]{0.24,0.48,0.48}{##1}}}
\@namedef{PY@tok@c1}{\let\PY@it=\textit\def\PY@tc##1{\textcolor[rgb]{0.24,0.48,0.48}{##1}}}
\@namedef{PY@tok@cs}{\let\PY@it=\textit\def\PY@tc##1{\textcolor[rgb]{0.24,0.48,0.48}{##1}}}

\def\PYZbs{\char`\\}
\def\PYZus{\char`\_}
\def\PYZob{\char`\{}
\def\PYZcb{\char`\}}
\def\PYZca{\char`\^}
\def\PYZam{\char`\&}
\def\PYZlt{\char`\<}
\def\PYZgt{\char`\>}
\def\PYZsh{\char`\#}
\def\PYZpc{\char`\%}
\def\PYZdl{\char`\$}
\def\PYZhy{\char`\-}
\def\PYZsq{\char`\'}
\def\PYZdq{\char`\"}
\def\PYZti{\char`\~}
% for compatibility with earlier versions
\def\PYZat{@}
\def\PYZlb{[}
\def\PYZrb{]}
\makeatother


    % For linebreaks inside Verbatim environment from package fancyvrb.
    \makeatletter
        \newbox\Wrappedcontinuationbox
        \newbox\Wrappedvisiblespacebox
        \newcommand*\Wrappedvisiblespace {\textcolor{red}{\textvisiblespace}}
        \newcommand*\Wrappedcontinuationsymbol {\textcolor{red}{\llap{\tiny$\m@th\hookrightarrow$}}}
        \newcommand*\Wrappedcontinuationindent {3ex }
        \newcommand*\Wrappedafterbreak {\kern\Wrappedcontinuationindent\copy\Wrappedcontinuationbox}
        % Take advantage of the already applied Pygments mark-up to insert
        % potential linebreaks for TeX processing.
        %        {, <, #, %, $, ' and ": go to next line.
        %        _, }, ^, &, >, - and ~: stay at end of broken line.
        % Use of \textquotesingle for straight quote.
        \newcommand*\Wrappedbreaksatspecials {%
            \def\PYGZus{\discretionary{\char`\_}{\Wrappedafterbreak}{\char`\_}}%
            \def\PYGZob{\discretionary{}{\Wrappedafterbreak\char`\{}{\char`\{}}%
            \def\PYGZcb{\discretionary{\char`\}}{\Wrappedafterbreak}{\char`\}}}%
            \def\PYGZca{\discretionary{\char`\^}{\Wrappedafterbreak}{\char`\^}}%
            \def\PYGZam{\discretionary{\char`\&}{\Wrappedafterbreak}{\char`\&}}%
            \def\PYGZlt{\discretionary{}{\Wrappedafterbreak\char`\<}{\char`\<}}%
            \def\PYGZgt{\discretionary{\char`\>}{\Wrappedafterbreak}{\char`\>}}%
            \def\PYGZsh{\discretionary{}{\Wrappedafterbreak\char`\#}{\char`\#}}%
            \def\PYGZpc{\discretionary{}{\Wrappedafterbreak\char`\%}{\char`\%}}%
            \def\PYGZdl{\discretionary{}{\Wrappedafterbreak\char`\$}{\char`\$}}%
            \def\PYGZhy{\discretionary{\char`\-}{\Wrappedafterbreak}{\char`\-}}%
            \def\PYGZsq{\discretionary{}{\Wrappedafterbreak\textquotesingle}{\textquotesingle}}%
            \def\PYGZdq{\discretionary{}{\Wrappedafterbreak\char`\"}{\char`\"}}%
            \def\PYGZti{\discretionary{\char`\~}{\Wrappedafterbreak}{\char`\~}}%
        }
        % Some characters . , ; ? ! / are not pygmentized.
        % This macro makes them "active" and they will insert potential linebreaks
        \newcommand*\Wrappedbreaksatpunct {%
            \lccode`\~`\.\lowercase{\def~}{\discretionary{\hbox{\char`\.}}{\Wrappedafterbreak}{\hbox{\char`\.}}}%
            \lccode`\~`\,\lowercase{\def~}{\discretionary{\hbox{\char`\,}}{\Wrappedafterbreak}{\hbox{\char`\,}}}%
            \lccode`\~`\;\lowercase{\def~}{\discretionary{\hbox{\char`\;}}{\Wrappedafterbreak}{\hbox{\char`\;}}}%
            \lccode`\~`\:\lowercase{\def~}{\discretionary{\hbox{\char`\:}}{\Wrappedafterbreak}{\hbox{\char`\:}}}%
            \lccode`\~`\?\lowercase{\def~}{\discretionary{\hbox{\char`\?}}{\Wrappedafterbreak}{\hbox{\char`\?}}}%
            \lccode`\~`\!\lowercase{\def~}{\discretionary{\hbox{\char`\!}}{\Wrappedafterbreak}{\hbox{\char`\!}}}%
            \lccode`\~`\/\lowercase{\def~}{\discretionary{\hbox{\char`\/}}{\Wrappedafterbreak}{\hbox{\char`\/}}}%
            \catcode`\.\active
            \catcode`\,\active
            \catcode`\;\active
            \catcode`\:\active
            \catcode`\?\active
            \catcode`\!\active
            \catcode`\/\active
            \lccode`\~`\~
        }
    \makeatother

    \let\OriginalVerbatim=\Verbatim
    \makeatletter
    \renewcommand{\Verbatim}[1][1]{%
        %\parskip\z@skip
        \sbox\Wrappedcontinuationbox {\Wrappedcontinuationsymbol}%
        \sbox\Wrappedvisiblespacebox {\FV@SetupFont\Wrappedvisiblespace}%
        \def\FancyVerbFormatLine ##1{\hsize\linewidth
            \vtop{\raggedright\hyphenpenalty\z@\exhyphenpenalty\z@
                \doublehyphendemerits\z@\finalhyphendemerits\z@
                \strut ##1\strut}%
        }%
        % If the linebreak is at a space, the latter will be displayed as visible
        % space at end of first line, and a continuation symbol starts next line.
        % Stretch/shrink are however usually zero for typewriter font.
        \def\FV@Space {%
            \nobreak\hskip\z@ plus\fontdimen3\font minus\fontdimen4\font
            \discretionary{\copy\Wrappedvisiblespacebox}{\Wrappedafterbreak}
            {\kern\fontdimen2\font}%
        }%

        % Allow breaks at special characters using \PYG... macros.
        \Wrappedbreaksatspecials
        % Breaks at punctuation characters . , ; ? ! and / need catcode=\active
        \OriginalVerbatim[#1,codes*=\Wrappedbreaksatpunct]%
    }
    \makeatother

    % Exact colors from NB
    \definecolor{incolor}{HTML}{303F9F}
    \definecolor{outcolor}{HTML}{D84315}
    \definecolor{cellborder}{HTML}{CFCFCF}
    \definecolor{cellbackground}{HTML}{F7F7F7}

    % prompt
    \makeatletter
    \newcommand{\boxspacing}{\kern\kvtcb@left@rule\kern\kvtcb@boxsep}
    \makeatother
    \newcommand{\prompt}[4]{
        {\ttfamily\llap{{\color{#2}[#3]:\hspace{3pt}#4}}\vspace{-\baselineskip}}
    }
    

    
    % Prevent overflowing lines due to hard-to-break entities
    \sloppy
    % Setup hyperref package
    \hypersetup{
      breaklinks=true,  % so long urls are correctly broken across lines
      colorlinks=true,
      urlcolor=urlcolor,
      linkcolor=linkcolor,
      citecolor=citecolor,
      }
    % Slightly bigger margins than the latex defaults
    
    \geometry{verbose,tmargin=1in,bmargin=1in,lmargin=1in,rmargin=1in}
    
    

\begin{document}
    
    \maketitle
    
    

    
    \hypertarget{linear-regression}{%
\section{Linear Regression}\label{linear-regression}}

    \hypertarget{importing-packages}{%
\subsection{Importing packages}\label{importing-packages}}

We import our standard libraries at this top level.

    \begin{tcolorbox}[breakable, size=fbox, boxrule=1pt, pad at break*=1mm,colback=cellbackground, colframe=cellborder]
\prompt{In}{incolor}{1}{\boxspacing}
\begin{Verbatim}[commandchars=\\\{\}]
\PY{k+kn}{import}\PY{+w}{ }\PY{n+nn}{numpy}\PY{+w}{ }\PY{k}{as}\PY{+w}{ }\PY{n+nn}{np}
\PY{k+kn}{import}\PY{+w}{ }\PY{n+nn}{pandas}\PY{+w}{ }\PY{k}{as}\PY{+w}{ }\PY{n+nn}{pd}
\PY{k+kn}{from}\PY{+w}{ }\PY{n+nn}{matplotlib}\PY{n+nn}{.}\PY{n+nn}{pyplot}\PY{+w}{ }\PY{k+kn}{import} \PY{n}{subplots}
\end{Verbatim}
\end{tcolorbox}

    \hypertarget{new-imports}{%
\subsubsection{New imports}\label{new-imports}}

Throughout this lab we will introduce new functions and libraries.
However, we will import them here to emphasize these are the new code
objects in this lab. Keeping imports near the top of a notebook makes
the code more readable, since scanning the first few lines tells us what
libraries are used.

    \begin{tcolorbox}[breakable, size=fbox, boxrule=1pt, pad at break*=1mm,colback=cellbackground, colframe=cellborder]
\prompt{In}{incolor}{2}{\boxspacing}
\begin{Verbatim}[commandchars=\\\{\}]
\PY{k+kn}{import}\PY{+w}{ }\PY{n+nn}{statsmodels}\PY{n+nn}{.}\PY{n+nn}{api}\PY{+w}{ }\PY{k}{as}\PY{+w}{ }\PY{n+nn}{sm}
\end{Verbatim}
\end{tcolorbox}

    We will provide relevant details about the functions below as they are
needed.

Besides importing whole modules, it is also possible to import only a
few items from a given module. This will help keep the \emph{namespace}
clean. We will use a few specific objects from the \texttt{statsmodels}
package which we import here.

    \begin{tcolorbox}[breakable, size=fbox, boxrule=1pt, pad at break*=1mm,colback=cellbackground, colframe=cellborder]
\prompt{In}{incolor}{3}{\boxspacing}
\begin{Verbatim}[commandchars=\\\{\}]
\PY{k+kn}{from}\PY{+w}{ }\PY{n+nn}{statsmodels}\PY{n+nn}{.}\PY{n+nn}{stats}\PY{n+nn}{.}\PY{n+nn}{outliers\PYZus{}influence} \PYZbs{}
     \PY{k+kn}{import}\PY{+w}{ }\PY{n+nn}{variance\PYZus{}inflation\PYZus{}factor}\PY{+w}{ }\PY{k}{as}\PY{+w}{ }\PY{n+nn}{VIF}
\PY{k+kn}{from}\PY{+w}{ }\PY{n+nn}{statsmodels}\PY{n+nn}{.}\PY{n+nn}{stats}\PY{n+nn}{.}\PY{n+nn}{anova}\PY{+w}{ }\PY{k+kn}{import} \PY{n}{anova\PYZus{}lm}
\end{Verbatim}
\end{tcolorbox}

    As one of the import statements above is quite a long line, we inserted
a line break \texttt{\textbackslash{}} to ease readability.

We will also use some functions written for the labs in this book in the
\texttt{ISLP} package.

    \begin{tcolorbox}[breakable, size=fbox, boxrule=1pt, pad at break*=1mm,colback=cellbackground, colframe=cellborder]
\prompt{In}{incolor}{4}{\boxspacing}
\begin{Verbatim}[commandchars=\\\{\}]
\PY{k+kn}{from}\PY{+w}{ }\PY{n+nn}{ISLP}\PY{+w}{ }\PY{k+kn}{import} \PY{n}{load\PYZus{}data}
\PY{k+kn}{from}\PY{+w}{ }\PY{n+nn}{ISLP}\PY{n+nn}{.}\PY{n+nn}{models}\PY{+w}{ }\PY{k+kn}{import} \PY{p}{(}\PY{n}{ModelSpec} \PY{k}{as} \PY{n}{MS}\PY{p}{,}
                         \PY{n}{summarize}\PY{p}{,}
                         \PY{n}{poly}\PY{p}{)}
\end{Verbatim}
\end{tcolorbox}

    \hypertarget{inspecting-objects-and-namespaces}{%
\subsubsection{Inspecting Objects and
Namespaces}\label{inspecting-objects-and-namespaces}}

The function \texttt{dir()} provides a list of objects in a namespace.

    \begin{tcolorbox}[breakable, size=fbox, boxrule=1pt, pad at break*=1mm,colback=cellbackground, colframe=cellborder]
\prompt{In}{incolor}{5}{\boxspacing}
\begin{Verbatim}[commandchars=\\\{\}]
\PY{n+nb}{dir}\PY{p}{(}\PY{p}{)}
\end{Verbatim}
\end{tcolorbox}

            \begin{tcolorbox}[breakable, size=fbox, boxrule=.5pt, pad at break*=1mm, opacityfill=0]
\prompt{Out}{outcolor}{5}{\boxspacing}
\begin{Verbatim}[commandchars=\\\{\}]
['In',
 'MS',
 'Out',
 'VIF',
 '\_',
 '\_\_',
 '\_\_\_',
 '\_\_builtin\_\_',
 '\_\_builtins\_\_',
 '\_\_doc\_\_',
 '\_\_loader\_\_',
 '\_\_name\_\_',
 '\_\_package\_\_',
 '\_\_spec\_\_',
 '\_dh',
 '\_i',
 '\_i1',
 '\_i2',
 '\_i3',
 '\_i4',
 '\_i5',
 '\_ih',
 '\_ii',
 '\_iii',
 '\_oh',
 'anova\_lm',
 'exit',
 'get\_ipython',
 'load\_data',
 'np',
 'open',
 'pd',
 'poly',
 'quit',
 'sm',
 'subplots',
 'summarize']
\end{Verbatim}
\end{tcolorbox}
        
    This shows you everything that \texttt{Python} can find at the top
level. There are certain objects like \texttt{\_\_builtins\_\_} that
contain references to built-in functions like \texttt{print()}.

Every python object has its own notion of namespace, also accessible
with \texttt{dir()}. This will include both the attributes of the object
as well as any methods associated with it. For instance, we see
\texttt{\textquotesingle{}sum\textquotesingle{}} in the listing for an
array.

    \begin{tcolorbox}[breakable, size=fbox, boxrule=1pt, pad at break*=1mm,colback=cellbackground, colframe=cellborder]
\prompt{In}{incolor}{6}{\boxspacing}
\begin{Verbatim}[commandchars=\\\{\}]
\PY{n}{A} \PY{o}{=} \PY{n}{np}\PY{o}{.}\PY{n}{array}\PY{p}{(}\PY{p}{[}\PY{l+m+mi}{3}\PY{p}{,}\PY{l+m+mi}{5}\PY{p}{,}\PY{l+m+mi}{11}\PY{p}{]}\PY{p}{)}
\PY{n+nb}{dir}\PY{p}{(}\PY{n}{A}\PY{p}{)}
\end{Verbatim}
\end{tcolorbox}

            \begin{tcolorbox}[breakable, size=fbox, boxrule=.5pt, pad at break*=1mm, opacityfill=0]
\prompt{Out}{outcolor}{6}{\boxspacing}
\begin{Verbatim}[commandchars=\\\{\}]
['T',
 '\_\_abs\_\_',
 '\_\_add\_\_',
 '\_\_and\_\_',
 '\_\_array\_\_',
 '\_\_array\_finalize\_\_',
 '\_\_array\_function\_\_',
 '\_\_array\_interface\_\_',
 '\_\_array\_prepare\_\_',
 '\_\_array\_priority\_\_',
 '\_\_array\_struct\_\_',
 '\_\_array\_ufunc\_\_',
 '\_\_array\_wrap\_\_',
 '\_\_bool\_\_',
 '\_\_class\_\_',
 '\_\_class\_getitem\_\_',
 '\_\_complex\_\_',
 '\_\_contains\_\_',
 '\_\_copy\_\_',
 '\_\_deepcopy\_\_',
 '\_\_delattr\_\_',
 '\_\_delitem\_\_',
 '\_\_dir\_\_',
 '\_\_divmod\_\_',
 '\_\_dlpack\_\_',
 '\_\_dlpack\_device\_\_',
 '\_\_doc\_\_',
 '\_\_eq\_\_',
 '\_\_float\_\_',
 '\_\_floordiv\_\_',
 '\_\_format\_\_',
 '\_\_ge\_\_',
 '\_\_getattribute\_\_',
 '\_\_getitem\_\_',
 '\_\_getstate\_\_',
 '\_\_gt\_\_',
 '\_\_hash\_\_',
 '\_\_iadd\_\_',
 '\_\_iand\_\_',
 '\_\_ifloordiv\_\_',
 '\_\_ilshift\_\_',
 '\_\_imatmul\_\_',
 '\_\_imod\_\_',
 '\_\_imul\_\_',
 '\_\_index\_\_',
 '\_\_init\_\_',
 '\_\_init\_subclass\_\_',
 '\_\_int\_\_',
 '\_\_invert\_\_',
 '\_\_ior\_\_',
 '\_\_ipow\_\_',
 '\_\_irshift\_\_',
 '\_\_isub\_\_',
 '\_\_iter\_\_',
 '\_\_itruediv\_\_',
 '\_\_ixor\_\_',
 '\_\_le\_\_',
 '\_\_len\_\_',
 '\_\_lshift\_\_',
 '\_\_lt\_\_',
 '\_\_matmul\_\_',
 '\_\_mod\_\_',
 '\_\_mul\_\_',
 '\_\_ne\_\_',
 '\_\_neg\_\_',
 '\_\_new\_\_',
 '\_\_or\_\_',
 '\_\_pos\_\_',
 '\_\_pow\_\_',
 '\_\_radd\_\_',
 '\_\_rand\_\_',
 '\_\_rdivmod\_\_',
 '\_\_reduce\_\_',
 '\_\_reduce\_ex\_\_',
 '\_\_repr\_\_',
 '\_\_rfloordiv\_\_',
 '\_\_rlshift\_\_',
 '\_\_rmatmul\_\_',
 '\_\_rmod\_\_',
 '\_\_rmul\_\_',
 '\_\_ror\_\_',
 '\_\_rpow\_\_',
 '\_\_rrshift\_\_',
 '\_\_rshift\_\_',
 '\_\_rsub\_\_',
 '\_\_rtruediv\_\_',
 '\_\_rxor\_\_',
 '\_\_setattr\_\_',
 '\_\_setitem\_\_',
 '\_\_setstate\_\_',
 '\_\_sizeof\_\_',
 '\_\_str\_\_',
 '\_\_sub\_\_',
 '\_\_subclasshook\_\_',
 '\_\_truediv\_\_',
 '\_\_xor\_\_',
 'all',
 'any',
 'argmax',
 'argmin',
 'argpartition',
 'argsort',
 'astype',
 'base',
 'byteswap',
 'choose',
 'clip',
 'compress',
 'conj',
 'conjugate',
 'copy',
 'ctypes',
 'cumprod',
 'cumsum',
 'data',
 'diagonal',
 'dot',
 'dtype',
 'dump',
 'dumps',
 'fill',
 'flags',
 'flat',
 'flatten',
 'getfield',
 'imag',
 'item',
 'itemset',
 'itemsize',
 'max',
 'mean',
 'min',
 'nbytes',
 'ndim',
 'newbyteorder',
 'nonzero',
 'partition',
 'prod',
 'ptp',
 'put',
 'ravel',
 'real',
 'repeat',
 'reshape',
 'resize',
 'round',
 'searchsorted',
 'setfield',
 'setflags',
 'shape',
 'size',
 'sort',
 'squeeze',
 'std',
 'strides',
 'sum',
 'swapaxes',
 'take',
 'tobytes',
 'tofile',
 'tolist',
 'tostring',
 'trace',
 'transpose',
 'var',
 'view']
\end{Verbatim}
\end{tcolorbox}
        
    This indicates that the object \texttt{A.sum} exists. In this case it is
a method that can be used to compute the sum of the array \texttt{A} as
can be seen by typing \texttt{A.sum?}.

    \begin{tcolorbox}[breakable, size=fbox, boxrule=1pt, pad at break*=1mm,colback=cellbackground, colframe=cellborder]
\prompt{In}{incolor}{7}{\boxspacing}
\begin{Verbatim}[commandchars=\\\{\}]
\PY{n}{A}\PY{o}{.}\PY{n}{sum}\PY{p}{(}\PY{p}{)}
\end{Verbatim}
\end{tcolorbox}

            \begin{tcolorbox}[breakable, size=fbox, boxrule=.5pt, pad at break*=1mm, opacityfill=0]
\prompt{Out}{outcolor}{7}{\boxspacing}
\begin{Verbatim}[commandchars=\\\{\}]
19
\end{Verbatim}
\end{tcolorbox}
        
    

    \hypertarget{simple-linear-regression}{%
\subsection{Simple Linear Regression}\label{simple-linear-regression}}

In this section we will construct model matrices (also called design
matrices) using the \texttt{ModelSpec()} transform from
\texttt{ISLP.models}.

We will use the \texttt{Boston} housing data set, which is contained in
the \texttt{ISLP} package. The \texttt{Boston} dataset records
\texttt{medv} (median house value) for \(506\) neighborhoods around
Boston. We will build a regression model to predict \texttt{medv} using
\(13\) predictors such as \texttt{rm} (average number of rooms per
house), \texttt{age} (proportion of owner-occupied units built prior to
1940), and \texttt{lstat} (percent of households with low socioeconomic
status). We will use \texttt{statsmodels} for this task, a
\texttt{Python} package that implements several commonly used regression
methods.

We have included a simple loading function \texttt{load\_data()} in the
\texttt{ISLP} package:

    \begin{tcolorbox}[breakable, size=fbox, boxrule=1pt, pad at break*=1mm,colback=cellbackground, colframe=cellborder]
\prompt{In}{incolor}{8}{\boxspacing}
\begin{Verbatim}[commandchars=\\\{\}]
\PY{n}{Boston} \PY{o}{=} \PY{n}{load\PYZus{}data}\PY{p}{(}\PY{l+s+s2}{\PYZdq{}}\PY{l+s+s2}{Boston}\PY{l+s+s2}{\PYZdq{}}\PY{p}{)}
\PY{n}{Boston}\PY{o}{.}\PY{n}{columns}
\end{Verbatim}
\end{tcolorbox}

            \begin{tcolorbox}[breakable, size=fbox, boxrule=.5pt, pad at break*=1mm, opacityfill=0]
\prompt{Out}{outcolor}{8}{\boxspacing}
\begin{Verbatim}[commandchars=\\\{\}]
Index(['crim', 'zn', 'indus', 'chas', 'nox', 'rm', 'age', 'dis', 'rad', 'tax',
       'ptratio', 'lstat', 'medv'],
      dtype='object')
\end{Verbatim}
\end{tcolorbox}
        
    Type \texttt{Boston?} to find out more about these data.

We start by using the \texttt{sm.OLS()} function to fit a simple linear
regression model. Our response will be \texttt{medv} and \texttt{lstat}
will be the single predictor. For this model, we can create the model
matrix by hand.

    \begin{tcolorbox}[breakable, size=fbox, boxrule=1pt, pad at break*=1mm,colback=cellbackground, colframe=cellborder]
\prompt{In}{incolor}{9}{\boxspacing}
\begin{Verbatim}[commandchars=\\\{\}]
\PY{n}{X} \PY{o}{=} \PY{n}{pd}\PY{o}{.}\PY{n}{DataFrame}\PY{p}{(}\PY{p}{\PYZob{}}\PY{l+s+s1}{\PYZsq{}}\PY{l+s+s1}{intercept}\PY{l+s+s1}{\PYZsq{}}\PY{p}{:} \PY{n}{np}\PY{o}{.}\PY{n}{ones}\PY{p}{(}\PY{n}{Boston}\PY{o}{.}\PY{n}{shape}\PY{p}{[}\PY{l+m+mi}{0}\PY{p}{]}\PY{p}{)}\PY{p}{,}
                  \PY{l+s+s1}{\PYZsq{}}\PY{l+s+s1}{lstat}\PY{l+s+s1}{\PYZsq{}}\PY{p}{:} \PY{n}{Boston}\PY{p}{[}\PY{l+s+s1}{\PYZsq{}}\PY{l+s+s1}{lstat}\PY{l+s+s1}{\PYZsq{}}\PY{p}{]}\PY{p}{\PYZcb{}}\PY{p}{)}
\PY{n}{X}\PY{p}{[}\PY{p}{:}\PY{l+m+mi}{4}\PY{p}{]}
\end{Verbatim}
\end{tcolorbox}

            \begin{tcolorbox}[breakable, size=fbox, boxrule=.5pt, pad at break*=1mm, opacityfill=0]
\prompt{Out}{outcolor}{9}{\boxspacing}
\begin{Verbatim}[commandchars=\\\{\}]
   intercept  lstat
0        1.0   4.98
1        1.0   9.14
2        1.0   4.03
3        1.0   2.94
\end{Verbatim}
\end{tcolorbox}
        
    We extract the response, and fit the model.

    \begin{tcolorbox}[breakable, size=fbox, boxrule=1pt, pad at break*=1mm,colback=cellbackground, colframe=cellborder]
\prompt{In}{incolor}{10}{\boxspacing}
\begin{Verbatim}[commandchars=\\\{\}]
\PY{n}{y} \PY{o}{=} \PY{n}{Boston}\PY{p}{[}\PY{l+s+s1}{\PYZsq{}}\PY{l+s+s1}{medv}\PY{l+s+s1}{\PYZsq{}}\PY{p}{]}
\PY{n}{model} \PY{o}{=} \PY{n}{sm}\PY{o}{.}\PY{n}{OLS}\PY{p}{(}\PY{n}{y}\PY{p}{,} \PY{n}{X}\PY{p}{)}
\PY{n}{results} \PY{o}{=} \PY{n}{model}\PY{o}{.}\PY{n}{fit}\PY{p}{(}\PY{p}{)}
\end{Verbatim}
\end{tcolorbox}

    Note that \texttt{sm.OLS()} does not fit the model; it specifies the
model, and then \texttt{model.fit()} does the actual fitting.

Our \texttt{ISLP} function \texttt{summarize()} produces a simple table
of the parameter estimates, their standard errors, t-statistics and
p-values. The function takes a single argument, such as the object
\texttt{results} returned here by the \texttt{fit} method, and returns
such a summary.

    \begin{tcolorbox}[breakable, size=fbox, boxrule=1pt, pad at break*=1mm,colback=cellbackground, colframe=cellborder]
\prompt{In}{incolor}{11}{\boxspacing}
\begin{Verbatim}[commandchars=\\\{\}]
\PY{n}{summarize}\PY{p}{(}\PY{n}{results}\PY{p}{)}
\end{Verbatim}
\end{tcolorbox}

            \begin{tcolorbox}[breakable, size=fbox, boxrule=.5pt, pad at break*=1mm, opacityfill=0]
\prompt{Out}{outcolor}{11}{\boxspacing}
\begin{Verbatim}[commandchars=\\\{\}]
              coef  std err       t  P>|t|
intercept  34.5538    0.563  61.415    0.0
lstat      -0.9500    0.039 -24.528    0.0
\end{Verbatim}
\end{tcolorbox}
        
    Before we describe other methods for working with fitted models, we
outline a more useful and general framework for constructing a model
matrix\textasciitilde{}\texttt{X}. \#\#\# Using Transformations: Fit and
Transform Our model above has a single predictor, and constructing
\texttt{X} was straightforward. In practice we often fit models with
more than one predictor, typically selected from an array or data frame.
We may wish to introduce transformations to the variables before fitting
the model, specify interactions between variables, and expand some
particular variables into sets of variables (e.g.~polynomials). The
\texttt{sklearn} package has a particular notion for this type of task:
a \emph{transform}. A transform is an object that is created with some
parameters as arguments. The object has two main methods: \texttt{fit()}
and \texttt{transform()}.

We provide a general approach for specifying models and constructing the
model matrix through the transform \texttt{ModelSpec()} in the
\texttt{ISLP} library. \texttt{ModelSpec()} (renamed \texttt{MS()} in
the preamble) creates a transform object, and then a pair of methods
\texttt{transform()} and \texttt{fit()} are used to construct a
corresponding model matrix.

We first describe this process for our simple regression model using a
single predictor \texttt{lstat} in the \texttt{Boston} data frame, but
will use it repeatedly in more complex tasks in this and other labs in
this book. In our case the transform is created by the expression
\texttt{design\ =\ MS({[}\textquotesingle{}lstat\textquotesingle{}{]})}.

The \texttt{fit()} method takes the original array and may do some
initial computations on it, as specified in the transform object. For
example, it may compute means and standard deviations for centering and
scaling. The \texttt{transform()} method applies the fitted
transformation to the array of data, and produces the model matrix.

    \begin{tcolorbox}[breakable, size=fbox, boxrule=1pt, pad at break*=1mm,colback=cellbackground, colframe=cellborder]
\prompt{In}{incolor}{12}{\boxspacing}
\begin{Verbatim}[commandchars=\\\{\}]
\PY{n}{design} \PY{o}{=} \PY{n}{MS}\PY{p}{(}\PY{p}{[}\PY{l+s+s1}{\PYZsq{}}\PY{l+s+s1}{lstat}\PY{l+s+s1}{\PYZsq{}}\PY{p}{]}\PY{p}{)}
\PY{n}{design} \PY{o}{=} \PY{n}{design}\PY{o}{.}\PY{n}{fit}\PY{p}{(}\PY{n}{Boston}\PY{p}{)}
\PY{n}{X} \PY{o}{=} \PY{n}{design}\PY{o}{.}\PY{n}{transform}\PY{p}{(}\PY{n}{Boston}\PY{p}{)}
\PY{n}{X}\PY{p}{[}\PY{p}{:}\PY{l+m+mi}{4}\PY{p}{]}
\end{Verbatim}
\end{tcolorbox}

            \begin{tcolorbox}[breakable, size=fbox, boxrule=.5pt, pad at break*=1mm, opacityfill=0]
\prompt{Out}{outcolor}{12}{\boxspacing}
\begin{Verbatim}[commandchars=\\\{\}]
   intercept  lstat
0        1.0   4.98
1        1.0   9.14
2        1.0   4.03
3        1.0   2.94
\end{Verbatim}
\end{tcolorbox}
        
    In this simple case, the \texttt{fit()} method does very little; it
simply checks that the variable
\texttt{\textquotesingle{}lstat\textquotesingle{}} specified in
\texttt{design} exists in \texttt{Boston}. Then \texttt{transform()}
constructs the model matrix with two columns: an \texttt{intercept} and
the variable \texttt{lstat}.

These two operations can be combined with the \texttt{fit\_transform()}
method.

    \begin{tcolorbox}[breakable, size=fbox, boxrule=1pt, pad at break*=1mm,colback=cellbackground, colframe=cellborder]
\prompt{In}{incolor}{13}{\boxspacing}
\begin{Verbatim}[commandchars=\\\{\}]
\PY{n}{design} \PY{o}{=} \PY{n}{MS}\PY{p}{(}\PY{p}{[}\PY{l+s+s1}{\PYZsq{}}\PY{l+s+s1}{lstat}\PY{l+s+s1}{\PYZsq{}}\PY{p}{]}\PY{p}{)}
\PY{n}{X} \PY{o}{=} \PY{n}{design}\PY{o}{.}\PY{n}{fit\PYZus{}transform}\PY{p}{(}\PY{n}{Boston}\PY{p}{)}
\PY{n}{X}\PY{p}{[}\PY{p}{:}\PY{l+m+mi}{4}\PY{p}{]}
\end{Verbatim}
\end{tcolorbox}

            \begin{tcolorbox}[breakable, size=fbox, boxrule=.5pt, pad at break*=1mm, opacityfill=0]
\prompt{Out}{outcolor}{13}{\boxspacing}
\begin{Verbatim}[commandchars=\\\{\}]
   intercept  lstat
0        1.0   4.98
1        1.0   9.14
2        1.0   4.03
3        1.0   2.94
\end{Verbatim}
\end{tcolorbox}
        
    Note that, as in the previous code chunk when the two steps were done
separately, the \texttt{design} object is changed as a result of the
\texttt{fit()} operation. The power of this pipeline will become clearer
when we fit more complex models that involve interactions and
transformations.

    Let's return to our fitted regression model. The object \texttt{results}
has several methods that can be used for inference. We already presented
a function \texttt{summarize()} for showing the essentials of the fit.
For a full and somewhat exhaustive summary of the fit, we can use the
\texttt{summary()} method.

    \begin{tcolorbox}[breakable, size=fbox, boxrule=1pt, pad at break*=1mm,colback=cellbackground, colframe=cellborder]
\prompt{In}{incolor}{14}{\boxspacing}
\begin{Verbatim}[commandchars=\\\{\}]
\PY{n}{results}\PY{o}{.}\PY{n}{summary}\PY{p}{(}\PY{p}{)}
\end{Verbatim}
\end{tcolorbox}
 
            
\prompt{Out}{outcolor}{14}{}
    
    \begin{center}
\begin{tabular}{lclc}
\toprule
\textbf{Dep. Variable:}    &       medv       & \textbf{  R-squared:         } &     0.544   \\
\textbf{Model:}            &       OLS        & \textbf{  Adj. R-squared:    } &     0.543   \\
\textbf{Method:}           &  Least Squares   & \textbf{  F-statistic:       } &     601.6   \\
\textbf{Date:}             & Thu, 03 Apr 2025 & \textbf{  Prob (F-statistic):} &  5.08e-88   \\
\textbf{Time:}             &     15:32:50     & \textbf{  Log-Likelihood:    } &   -1641.5   \\
\textbf{No. Observations:} &         506      & \textbf{  AIC:               } &     3287.   \\
\textbf{Df Residuals:}     &         504      & \textbf{  BIC:               } &     3295.   \\
\textbf{Df Model:}         &           1      & \textbf{                     } &             \\
\textbf{Covariance Type:}  &    nonrobust     & \textbf{                     } &             \\
\bottomrule
\end{tabular}
\begin{tabular}{lcccccc}
                   & \textbf{coef} & \textbf{std err} & \textbf{t} & \textbf{P$> |$t$|$} & \textbf{[0.025} & \textbf{0.975]}  \\
\midrule
\textbf{intercept} &      34.5538  &        0.563     &    61.415  &         0.000        &       33.448    &       35.659     \\
\textbf{lstat}     &      -0.9500  &        0.039     &   -24.528  &         0.000        &       -1.026    &       -0.874     \\
\bottomrule
\end{tabular}
\begin{tabular}{lclc}
\textbf{Omnibus:}       & 137.043 & \textbf{  Durbin-Watson:     } &    0.892  \\
\textbf{Prob(Omnibus):} &   0.000 & \textbf{  Jarque-Bera (JB):  } &  291.373  \\
\textbf{Skew:}          &   1.453 & \textbf{  Prob(JB):          } & 5.36e-64  \\
\textbf{Kurtosis:}      &   5.319 & \textbf{  Cond. No.          } &     29.7  \\
\bottomrule
\end{tabular}
%\caption{OLS Regression Results}
\end{center}

Notes: \newline
 [1] Standard Errors assume that the covariance matrix of the errors is correctly specified.

    

    The fitted coefficients can also be retrieved as the \texttt{params}
attribute of \texttt{results}.

    \begin{tcolorbox}[breakable, size=fbox, boxrule=1pt, pad at break*=1mm,colback=cellbackground, colframe=cellborder]
\prompt{In}{incolor}{15}{\boxspacing}
\begin{Verbatim}[commandchars=\\\{\}]
\PY{n}{results}\PY{o}{.}\PY{n}{params}
\end{Verbatim}
\end{tcolorbox}

            \begin{tcolorbox}[breakable, size=fbox, boxrule=.5pt, pad at break*=1mm, opacityfill=0]
\prompt{Out}{outcolor}{15}{\boxspacing}
\begin{Verbatim}[commandchars=\\\{\}]
intercept    34.553841
lstat        -0.950049
dtype: float64
\end{Verbatim}
\end{tcolorbox}
        
    The \texttt{get\_prediction()} method can be used to obtain predictions,
and produce confidence intervals and prediction intervals for the
prediction of \texttt{medv} for given values of \texttt{lstat}.

We first create a new data frame, in this case containing only the
variable \texttt{lstat}, with the values for this variable at which we
wish to make predictions. We then use the \texttt{transform()} method of
\texttt{design} to create the corresponding model matrix.

    \begin{tcolorbox}[breakable, size=fbox, boxrule=1pt, pad at break*=1mm,colback=cellbackground, colframe=cellborder]
\prompt{In}{incolor}{16}{\boxspacing}
\begin{Verbatim}[commandchars=\\\{\}]
\PY{n}{new\PYZus{}df} \PY{o}{=} \PY{n}{pd}\PY{o}{.}\PY{n}{DataFrame}\PY{p}{(}\PY{p}{\PYZob{}}\PY{l+s+s1}{\PYZsq{}}\PY{l+s+s1}{lstat}\PY{l+s+s1}{\PYZsq{}}\PY{p}{:}\PY{p}{[}\PY{l+m+mi}{5}\PY{p}{,} \PY{l+m+mi}{10}\PY{p}{,} \PY{l+m+mi}{15}\PY{p}{]}\PY{p}{\PYZcb{}}\PY{p}{)}
\PY{n}{newX} \PY{o}{=} \PY{n}{design}\PY{o}{.}\PY{n}{transform}\PY{p}{(}\PY{n}{new\PYZus{}df}\PY{p}{)}
\PY{n}{newX}
\end{Verbatim}
\end{tcolorbox}

            \begin{tcolorbox}[breakable, size=fbox, boxrule=.5pt, pad at break*=1mm, opacityfill=0]
\prompt{Out}{outcolor}{16}{\boxspacing}
\begin{Verbatim}[commandchars=\\\{\}]
   intercept  lstat
0        1.0      5
1        1.0     10
2        1.0     15
\end{Verbatim}
\end{tcolorbox}
        
    Next we compute the predictions at \texttt{newX}, and view them by
extracting the \texttt{predicted\_mean} attribute.

    \begin{tcolorbox}[breakable, size=fbox, boxrule=1pt, pad at break*=1mm,colback=cellbackground, colframe=cellborder]
\prompt{In}{incolor}{17}{\boxspacing}
\begin{Verbatim}[commandchars=\\\{\}]
\PY{n}{new\PYZus{}predictions} \PY{o}{=} \PY{n}{results}\PY{o}{.}\PY{n}{get\PYZus{}prediction}\PY{p}{(}\PY{n}{newX}\PY{p}{)}\PY{p}{;}
\PY{n}{new\PYZus{}predictions}\PY{o}{.}\PY{n}{predicted\PYZus{}mean}
\end{Verbatim}
\end{tcolorbox}

            \begin{tcolorbox}[breakable, size=fbox, boxrule=.5pt, pad at break*=1mm, opacityfill=0]
\prompt{Out}{outcolor}{17}{\boxspacing}
\begin{Verbatim}[commandchars=\\\{\}]
array([29.80359411, 25.05334734, 20.30310057])
\end{Verbatim}
\end{tcolorbox}
        
    We can produce confidence intervals for the predicted values.

    \begin{tcolorbox}[breakable, size=fbox, boxrule=1pt, pad at break*=1mm,colback=cellbackground, colframe=cellborder]
\prompt{In}{incolor}{18}{\boxspacing}
\begin{Verbatim}[commandchars=\\\{\}]
\PY{n}{new\PYZus{}predictions}\PY{o}{.}\PY{n}{conf\PYZus{}int}\PY{p}{(}\PY{n}{alpha}\PY{o}{=}\PY{l+m+mf}{0.05}\PY{p}{)}
\end{Verbatim}
\end{tcolorbox}

            \begin{tcolorbox}[breakable, size=fbox, boxrule=.5pt, pad at break*=1mm, opacityfill=0]
\prompt{Out}{outcolor}{18}{\boxspacing}
\begin{Verbatim}[commandchars=\\\{\}]
array([[29.00741194, 30.59977628],
       [24.47413202, 25.63256267],
       [19.73158815, 20.87461299]])
\end{Verbatim}
\end{tcolorbox}
        
    Prediction intervals are computed by setting \texttt{obs=True}:

    \begin{tcolorbox}[breakable, size=fbox, boxrule=1pt, pad at break*=1mm,colback=cellbackground, colframe=cellborder]
\prompt{In}{incolor}{19}{\boxspacing}
\begin{Verbatim}[commandchars=\\\{\}]
\PY{n}{new\PYZus{}predictions}\PY{o}{.}\PY{n}{conf\PYZus{}int}\PY{p}{(}\PY{n}{obs}\PY{o}{=}\PY{k+kc}{True}\PY{p}{,} \PY{n}{alpha}\PY{o}{=}\PY{l+m+mf}{0.05}\PY{p}{)}
\end{Verbatim}
\end{tcolorbox}

            \begin{tcolorbox}[breakable, size=fbox, boxrule=.5pt, pad at break*=1mm, opacityfill=0]
\prompt{Out}{outcolor}{19}{\boxspacing}
\begin{Verbatim}[commandchars=\\\{\}]
array([[17.56567478, 42.04151344],
       [12.82762635, 37.27906833],
       [ 8.0777421 , 32.52845905]])
\end{Verbatim}
\end{tcolorbox}
        
    For instance, the 95\% confidence interval associated with an
\texttt{lstat} value of 10 is (24.47, 25.63), and the 95\% prediction
interval is (12.82, 37.28). As expected, the confidence and prediction
intervals are centered around the same point (a predicted value of 25.05
for \texttt{medv} when \texttt{lstat} equals 10), but the latter are
substantially wider.

Next we will plot \texttt{medv} and \texttt{lstat} using
\texttt{DataFrame.plot.scatter()} and wish to add the regression line to
the resulting plot.

    \hypertarget{defining-functions}{%
\subsubsection{Defining Functions}\label{defining-functions}}

While there is a function within the \texttt{ISLP} package that adds a
line to an existing plot, we take this opportunity to define our first
function to do so.

    \begin{tcolorbox}[breakable, size=fbox, boxrule=1pt, pad at break*=1mm,colback=cellbackground, colframe=cellborder]
\prompt{In}{incolor}{20}{\boxspacing}
\begin{Verbatim}[commandchars=\\\{\}]
\PY{k}{def}\PY{+w}{ }\PY{n+nf}{abline}\PY{p}{(}\PY{n}{ax}\PY{p}{,} \PY{n}{b}\PY{p}{,} \PY{n}{m}\PY{p}{)}\PY{p}{:}
    \PY{l+s+s2}{\PYZdq{}}\PY{l+s+s2}{Add a line with slope m and intercept b to ax}\PY{l+s+s2}{\PYZdq{}}
    \PY{n}{xlim} \PY{o}{=} \PY{n}{ax}\PY{o}{.}\PY{n}{get\PYZus{}xlim}\PY{p}{(}\PY{p}{)}
    \PY{n}{ylim} \PY{o}{=} \PY{p}{[}\PY{n}{m} \PY{o}{*} \PY{n}{xlim}\PY{p}{[}\PY{l+m+mi}{0}\PY{p}{]} \PY{o}{+} \PY{n}{b}\PY{p}{,} \PY{n}{m} \PY{o}{*} \PY{n}{xlim}\PY{p}{[}\PY{l+m+mi}{1}\PY{p}{]} \PY{o}{+} \PY{n}{b}\PY{p}{]}
    \PY{n}{ax}\PY{o}{.}\PY{n}{plot}\PY{p}{(}\PY{n}{xlim}\PY{p}{,} \PY{n}{ylim}\PY{p}{)}
\end{Verbatim}
\end{tcolorbox}

    A few things are illustrated above. First we see the syntax for defining
a function: \texttt{def\ funcname(...)}. The function has arguments
\texttt{ax,\ b,\ m} where \texttt{ax} is an axis object for an existing
plot, \texttt{b} is the intercept and \texttt{m} is the slope of the
desired line. Other plotting options can be passed on to
\texttt{ax.plot} by including additional optional arguments as follows:

    \begin{tcolorbox}[breakable, size=fbox, boxrule=1pt, pad at break*=1mm,colback=cellbackground, colframe=cellborder]
\prompt{In}{incolor}{21}{\boxspacing}
\begin{Verbatim}[commandchars=\\\{\}]
\PY{k}{def}\PY{+w}{ }\PY{n+nf}{abline}\PY{p}{(}\PY{n}{ax}\PY{p}{,} \PY{n}{b}\PY{p}{,} \PY{n}{m}\PY{p}{,} \PY{o}{*}\PY{n}{args}\PY{p}{,} \PY{o}{*}\PY{o}{*}\PY{n}{kwargs}\PY{p}{)}\PY{p}{:}
    \PY{l+s+s2}{\PYZdq{}}\PY{l+s+s2}{Add a line with slope m and intercept b to ax}\PY{l+s+s2}{\PYZdq{}}
    \PY{n}{xlim} \PY{o}{=} \PY{n}{ax}\PY{o}{.}\PY{n}{get\PYZus{}xlim}\PY{p}{(}\PY{p}{)}
    \PY{n}{ylim} \PY{o}{=} \PY{p}{[}\PY{n}{m} \PY{o}{*} \PY{n}{xlim}\PY{p}{[}\PY{l+m+mi}{0}\PY{p}{]} \PY{o}{+} \PY{n}{b}\PY{p}{,} \PY{n}{m} \PY{o}{*} \PY{n}{xlim}\PY{p}{[}\PY{l+m+mi}{1}\PY{p}{]} \PY{o}{+} \PY{n}{b}\PY{p}{]}
    \PY{n}{ax}\PY{o}{.}\PY{n}{plot}\PY{p}{(}\PY{n}{xlim}\PY{p}{,} \PY{n}{ylim}\PY{p}{,} \PY{o}{*}\PY{n}{args}\PY{p}{,} \PY{o}{*}\PY{o}{*}\PY{n}{kwargs}\PY{p}{)}
\end{Verbatim}
\end{tcolorbox}

    The addition of \texttt{*args} allows any number of non-named arguments
to \texttt{abline}, while \texttt{**kwargs} allows any number of named
arguments (such as \texttt{linewidth=3}) to \texttt{abline}. In our
function, we pass these arguments verbatim to \texttt{ax.plot} above.
Readers interested in learning more about functions are referred to the
section on defining functions in
\href{https://docs.python.org/3/tutorial/controlflow.html\#defining-functions}{docs.python.org/tutorial}.

Let's use our new function to add this regression line to a plot of
\texttt{medv} vs.~\texttt{lstat}.

    \begin{tcolorbox}[breakable, size=fbox, boxrule=1pt, pad at break*=1mm,colback=cellbackground, colframe=cellborder]
\prompt{In}{incolor}{22}{\boxspacing}
\begin{Verbatim}[commandchars=\\\{\}]
\PY{n}{ax} \PY{o}{=} \PY{n}{Boston}\PY{o}{.}\PY{n}{plot}\PY{o}{.}\PY{n}{scatter}\PY{p}{(}\PY{l+s+s1}{\PYZsq{}}\PY{l+s+s1}{lstat}\PY{l+s+s1}{\PYZsq{}}\PY{p}{,} \PY{l+s+s1}{\PYZsq{}}\PY{l+s+s1}{medv}\PY{l+s+s1}{\PYZsq{}}\PY{p}{)}
\PY{n}{abline}\PY{p}{(}\PY{n}{ax}\PY{p}{,}
       \PY{n}{results}\PY{o}{.}\PY{n}{params}\PY{p}{[}\PY{l+m+mi}{0}\PY{p}{]}\PY{p}{,}
       \PY{n}{results}\PY{o}{.}\PY{n}{params}\PY{p}{[}\PY{l+m+mi}{1}\PY{p}{]}\PY{p}{,}
       \PY{l+s+s1}{\PYZsq{}}\PY{l+s+s1}{r\PYZhy{}\PYZhy{}}\PY{l+s+s1}{\PYZsq{}}\PY{p}{,}
       \PY{n}{linewidth}\PY{o}{=}\PY{l+m+mi}{3}\PY{p}{)}
\end{Verbatim}
\end{tcolorbox}

    \begin{Verbatim}[commandchars=\\\{\}]
/var/folders/16/8y65\_zv174qgdp4ktlmpv12h0000gq/T/ipykernel\_70382/1591428221.py:3
: FutureWarning: Series.\_\_getitem\_\_ treating keys as positions is deprecated. In
a future version, integer keys will always be treated as labels (consistent with
DataFrame behavior). To access a value by position, use `ser.iloc[pos]`
  results.params[0],
/var/folders/16/8y65\_zv174qgdp4ktlmpv12h0000gq/T/ipykernel\_70382/1591428221.py:4
: FutureWarning: Series.\_\_getitem\_\_ treating keys as positions is deprecated. In
a future version, integer keys will always be treated as labels (consistent with
DataFrame behavior). To access a value by position, use `ser.iloc[pos]`
  results.params[1],
    \end{Verbatim}

    \begin{center}
    \adjustimage{max size={0.9\linewidth}{0.9\paperheight}}{artifacts/latex/Ch03-linreg-lab_files/artifacts/latex/Ch03-linreg-lab_47_1.png}
    \end{center}
    { \hspace*{\fill} \\}
    
    Thus, the final call to \texttt{ax.plot()} is
\texttt{ax.plot(xlim,\ ylim,\ \textquotesingle{}r-\/-\textquotesingle{},\ linewidth=3)}.
We have used the argument
\texttt{\textquotesingle{}r-\/-\textquotesingle{}} to produce a red
dashed line, and added an argument to make it of width 3. There is some
evidence for non-linearity in the relationship between \texttt{lstat}
and \texttt{medv}. We will explore this issue later in this lab.

As mentioned above, there is an existing function to add a line to a
plot --- \texttt{ax.axline()} --- but knowing how to write such
functions empowers us to create more expressive displays.

    Next we examine some diagnostic plots, several of which were discussed
in Section 3.3.3. We can find the fitted values and residuals of the fit
as attributes of the \texttt{results} object. Various influence measures
describing the regression model are computed with the
\texttt{get\_influence()} method. As we will not use the \texttt{fig}
component returned as the first value from \texttt{subplots()}, we
simply capture the second returned value in \texttt{ax} below.

    \begin{tcolorbox}[breakable, size=fbox, boxrule=1pt, pad at break*=1mm,colback=cellbackground, colframe=cellborder]
\prompt{In}{incolor}{23}{\boxspacing}
\begin{Verbatim}[commandchars=\\\{\}]
\PY{n}{ax} \PY{o}{=} \PY{n}{subplots}\PY{p}{(}\PY{n}{figsize}\PY{o}{=}\PY{p}{(}\PY{l+m+mi}{8}\PY{p}{,}\PY{l+m+mi}{8}\PY{p}{)}\PY{p}{)}\PY{p}{[}\PY{l+m+mi}{1}\PY{p}{]}
\PY{n}{ax}\PY{o}{.}\PY{n}{scatter}\PY{p}{(}\PY{n}{results}\PY{o}{.}\PY{n}{fittedvalues}\PY{p}{,} \PY{n}{results}\PY{o}{.}\PY{n}{resid}\PY{p}{)}
\PY{n}{ax}\PY{o}{.}\PY{n}{set\PYZus{}xlabel}\PY{p}{(}\PY{l+s+s1}{\PYZsq{}}\PY{l+s+s1}{Fitted value}\PY{l+s+s1}{\PYZsq{}}\PY{p}{)}
\PY{n}{ax}\PY{o}{.}\PY{n}{set\PYZus{}ylabel}\PY{p}{(}\PY{l+s+s1}{\PYZsq{}}\PY{l+s+s1}{Residual}\PY{l+s+s1}{\PYZsq{}}\PY{p}{)}
\PY{n}{ax}\PY{o}{.}\PY{n}{axhline}\PY{p}{(}\PY{l+m+mi}{0}\PY{p}{,} \PY{n}{c}\PY{o}{=}\PY{l+s+s1}{\PYZsq{}}\PY{l+s+s1}{k}\PY{l+s+s1}{\PYZsq{}}\PY{p}{,} \PY{n}{ls}\PY{o}{=}\PY{l+s+s1}{\PYZsq{}}\PY{l+s+s1}{\PYZhy{}\PYZhy{}}\PY{l+s+s1}{\PYZsq{}}\PY{p}{)}\PY{p}{;}
\end{Verbatim}
\end{tcolorbox}

    \begin{center}
    \adjustimage{max size={0.9\linewidth}{0.9\paperheight}}{artifacts/latex/Ch03-linreg-lab_files/artifacts/latex/Ch03-linreg-lab_50_0.png}
    \end{center}
    { \hspace*{\fill} \\}
    
    We add a horizontal line at 0 for reference using the
\texttt{ax.axhline()} method, indicating it should be black
(\texttt{c=\textquotesingle{}k\textquotesingle{}}) and have a dashed
linestyle (\texttt{ls=\textquotesingle{}-\/-\textquotesingle{}}).

On the basis of the residual plot, there is some evidence of
non-linearity. Leverage statistics can be computed for any number of
predictors using the \texttt{hat\_matrix\_diag} attribute of the value
returned by the \texttt{get\_influence()} method.

    \begin{tcolorbox}[breakable, size=fbox, boxrule=1pt, pad at break*=1mm,colback=cellbackground, colframe=cellborder]
\prompt{In}{incolor}{24}{\boxspacing}
\begin{Verbatim}[commandchars=\\\{\}]
\PY{n}{infl} \PY{o}{=} \PY{n}{results}\PY{o}{.}\PY{n}{get\PYZus{}influence}\PY{p}{(}\PY{p}{)}
\PY{n}{ax} \PY{o}{=} \PY{n}{subplots}\PY{p}{(}\PY{n}{figsize}\PY{o}{=}\PY{p}{(}\PY{l+m+mi}{8}\PY{p}{,}\PY{l+m+mi}{8}\PY{p}{)}\PY{p}{)}\PY{p}{[}\PY{l+m+mi}{1}\PY{p}{]}
\PY{n}{ax}\PY{o}{.}\PY{n}{scatter}\PY{p}{(}\PY{n}{np}\PY{o}{.}\PY{n}{arange}\PY{p}{(}\PY{n}{X}\PY{o}{.}\PY{n}{shape}\PY{p}{[}\PY{l+m+mi}{0}\PY{p}{]}\PY{p}{)}\PY{p}{,} \PY{n}{infl}\PY{o}{.}\PY{n}{hat\PYZus{}matrix\PYZus{}diag}\PY{p}{)}
\PY{n}{ax}\PY{o}{.}\PY{n}{set\PYZus{}xlabel}\PY{p}{(}\PY{l+s+s1}{\PYZsq{}}\PY{l+s+s1}{Index}\PY{l+s+s1}{\PYZsq{}}\PY{p}{)}
\PY{n}{ax}\PY{o}{.}\PY{n}{set\PYZus{}ylabel}\PY{p}{(}\PY{l+s+s1}{\PYZsq{}}\PY{l+s+s1}{Leverage}\PY{l+s+s1}{\PYZsq{}}\PY{p}{)}
\PY{n}{np}\PY{o}{.}\PY{n}{argmax}\PY{p}{(}\PY{n}{infl}\PY{o}{.}\PY{n}{hat\PYZus{}matrix\PYZus{}diag}\PY{p}{)}
\end{Verbatim}
\end{tcolorbox}

            \begin{tcolorbox}[breakable, size=fbox, boxrule=.5pt, pad at break*=1mm, opacityfill=0]
\prompt{Out}{outcolor}{24}{\boxspacing}
\begin{Verbatim}[commandchars=\\\{\}]
374
\end{Verbatim}
\end{tcolorbox}
        
    \begin{center}
    \adjustimage{max size={0.9\linewidth}{0.9\paperheight}}{artifacts/latex/Ch03-linreg-lab_files/artifacts/latex/Ch03-linreg-lab_52_1.png}
    \end{center}
    { \hspace*{\fill} \\}
    
    The \texttt{np.argmax()} function identifies the index of the largest
element of an array, optionally computed over an axis of the array. In
this case, we maximized over the entire array to determine which
observation has the largest leverage statistic.

    \hypertarget{multiple-linear-regression}{%
\subsection{Multiple Linear
Regression}\label{multiple-linear-regression}}

In order to fit a multiple linear regression model using least squares,
we again use the \texttt{ModelSpec()} transform to construct the
required model matrix and response. The arguments to
\texttt{ModelSpec()} can be quite general, but in this case a list of
column names suffice. We consider a fit here with the two variables
\texttt{lstat} and \texttt{age}.

    \begin{tcolorbox}[breakable, size=fbox, boxrule=1pt, pad at break*=1mm,colback=cellbackground, colframe=cellborder]
\prompt{In}{incolor}{25}{\boxspacing}
\begin{Verbatim}[commandchars=\\\{\}]
\PY{n}{X} \PY{o}{=} \PY{n}{MS}\PY{p}{(}\PY{p}{[}\PY{l+s+s1}{\PYZsq{}}\PY{l+s+s1}{lstat}\PY{l+s+s1}{\PYZsq{}}\PY{p}{,} \PY{l+s+s1}{\PYZsq{}}\PY{l+s+s1}{age}\PY{l+s+s1}{\PYZsq{}}\PY{p}{]}\PY{p}{)}\PY{o}{.}\PY{n}{fit\PYZus{}transform}\PY{p}{(}\PY{n}{Boston}\PY{p}{)}
\PY{n}{model1} \PY{o}{=} \PY{n}{sm}\PY{o}{.}\PY{n}{OLS}\PY{p}{(}\PY{n}{y}\PY{p}{,} \PY{n}{X}\PY{p}{)}
\PY{n}{results1} \PY{o}{=} \PY{n}{model1}\PY{o}{.}\PY{n}{fit}\PY{p}{(}\PY{p}{)}
\PY{n}{summarize}\PY{p}{(}\PY{n}{results1}\PY{p}{)}
\end{Verbatim}
\end{tcolorbox}

            \begin{tcolorbox}[breakable, size=fbox, boxrule=.5pt, pad at break*=1mm, opacityfill=0]
\prompt{Out}{outcolor}{25}{\boxspacing}
\begin{Verbatim}[commandchars=\\\{\}]
              coef  std err       t  P>|t|
intercept  33.2228    0.731  45.458  0.000
lstat      -1.0321    0.048 -21.416  0.000
age         0.0345    0.012   2.826  0.005
\end{Verbatim}
\end{tcolorbox}
        
    Notice how we have compacted the first line into a succinct expression
describing the construction of \texttt{X}.

The \texttt{Boston} data set contains 12 variables, and so it would be
cumbersome to have to type all of these in order to perform a regression
using all of the predictors. Instead, we can use the following
short-hand:\definelongblankMR{columns.drop()}{columns.slashslashdrop()}

    \begin{tcolorbox}[breakable, size=fbox, boxrule=1pt, pad at break*=1mm,colback=cellbackground, colframe=cellborder]
\prompt{In}{incolor}{26}{\boxspacing}
\begin{Verbatim}[commandchars=\\\{\}]
\PY{n}{terms} \PY{o}{=} \PY{n}{Boston}\PY{o}{.}\PY{n}{columns}\PY{o}{.}\PY{n}{drop}\PY{p}{(}\PY{l+s+s1}{\PYZsq{}}\PY{l+s+s1}{medv}\PY{l+s+s1}{\PYZsq{}}\PY{p}{)}
\PY{n}{terms}
\end{Verbatim}
\end{tcolorbox}

            \begin{tcolorbox}[breakable, size=fbox, boxrule=.5pt, pad at break*=1mm, opacityfill=0]
\prompt{Out}{outcolor}{26}{\boxspacing}
\begin{Verbatim}[commandchars=\\\{\}]
Index(['crim', 'zn', 'indus', 'chas', 'nox', 'rm', 'age', 'dis', 'rad', 'tax',
       'ptratio', 'lstat'],
      dtype='object')
\end{Verbatim}
\end{tcolorbox}
        
    We can now fit the model with all the variables in \texttt{terms} using
the same model matrix builder.

    \begin{tcolorbox}[breakable, size=fbox, boxrule=1pt, pad at break*=1mm,colback=cellbackground, colframe=cellborder]
\prompt{In}{incolor}{27}{\boxspacing}
\begin{Verbatim}[commandchars=\\\{\}]
\PY{n}{X} \PY{o}{=} \PY{n}{MS}\PY{p}{(}\PY{n}{terms}\PY{p}{)}\PY{o}{.}\PY{n}{fit\PYZus{}transform}\PY{p}{(}\PY{n}{Boston}\PY{p}{)}
\PY{n}{model} \PY{o}{=} \PY{n}{sm}\PY{o}{.}\PY{n}{OLS}\PY{p}{(}\PY{n}{y}\PY{p}{,} \PY{n}{X}\PY{p}{)}
\PY{n}{results} \PY{o}{=} \PY{n}{model}\PY{o}{.}\PY{n}{fit}\PY{p}{(}\PY{p}{)}
\PY{n}{summarize}\PY{p}{(}\PY{n}{results}\PY{p}{)}
\end{Verbatim}
\end{tcolorbox}

            \begin{tcolorbox}[breakable, size=fbox, boxrule=.5pt, pad at break*=1mm, opacityfill=0]
\prompt{Out}{outcolor}{27}{\boxspacing}
\begin{Verbatim}[commandchars=\\\{\}]
              coef  std err       t  P>|t|
intercept  41.6173    4.936   8.431  0.000
crim       -0.1214    0.033  -3.678  0.000
zn          0.0470    0.014   3.384  0.001
indus       0.0135    0.062   0.217  0.829
chas        2.8400    0.870   3.264  0.001
nox       -18.7580    3.851  -4.870  0.000
rm          3.6581    0.420   8.705  0.000
age         0.0036    0.013   0.271  0.787
dis        -1.4908    0.202  -7.394  0.000
rad         0.2894    0.067   4.325  0.000
tax        -0.0127    0.004  -3.337  0.001
ptratio    -0.9375    0.132  -7.091  0.000
lstat      -0.5520    0.051 -10.897  0.000
\end{Verbatim}
\end{tcolorbox}
        
    What if we would like to perform a regression using all of the variables
but one? For example, in the above regression output, \texttt{age} has a
high \(p\)-value. So we may wish to run a regression excluding this
predictor. The following syntax results in a regression using all
predictors except \texttt{age}.

    \begin{tcolorbox}[breakable, size=fbox, boxrule=1pt, pad at break*=1mm,colback=cellbackground, colframe=cellborder]
\prompt{In}{incolor}{28}{\boxspacing}
\begin{Verbatim}[commandchars=\\\{\}]
\PY{n}{minus\PYZus{}age} \PY{o}{=} \PY{n}{Boston}\PY{o}{.}\PY{n}{columns}\PY{o}{.}\PY{n}{drop}\PY{p}{(}\PY{p}{[}\PY{l+s+s1}{\PYZsq{}}\PY{l+s+s1}{medv}\PY{l+s+s1}{\PYZsq{}}\PY{p}{,} \PY{l+s+s1}{\PYZsq{}}\PY{l+s+s1}{age}\PY{l+s+s1}{\PYZsq{}}\PY{p}{]}\PY{p}{)} 
\PY{n}{Xma} \PY{o}{=} \PY{n}{MS}\PY{p}{(}\PY{n}{minus\PYZus{}age}\PY{p}{)}\PY{o}{.}\PY{n}{fit\PYZus{}transform}\PY{p}{(}\PY{n}{Boston}\PY{p}{)}
\PY{n}{model1} \PY{o}{=} \PY{n}{sm}\PY{o}{.}\PY{n}{OLS}\PY{p}{(}\PY{n}{y}\PY{p}{,} \PY{n}{Xma}\PY{p}{)}
\PY{n}{summarize}\PY{p}{(}\PY{n}{model1}\PY{o}{.}\PY{n}{fit}\PY{p}{(}\PY{p}{)}\PY{p}{)}
\end{Verbatim}
\end{tcolorbox}

            \begin{tcolorbox}[breakable, size=fbox, boxrule=.5pt, pad at break*=1mm, opacityfill=0]
\prompt{Out}{outcolor}{28}{\boxspacing}
\begin{Verbatim}[commandchars=\\\{\}]
              coef  std err       t  P>|t|
intercept  41.5251    4.920   8.441  0.000
crim       -0.1214    0.033  -3.683  0.000
zn          0.0465    0.014   3.379  0.001
indus       0.0135    0.062   0.217  0.829
chas        2.8528    0.868   3.287  0.001
nox       -18.4851    3.714  -4.978  0.000
rm          3.6811    0.411   8.951  0.000
dis        -1.5068    0.193  -7.825  0.000
rad         0.2879    0.067   4.322  0.000
tax        -0.0127    0.004  -3.333  0.001
ptratio    -0.9346    0.132  -7.099  0.000
lstat      -0.5474    0.048 -11.483  0.000
\end{Verbatim}
\end{tcolorbox}
        
    \hypertarget{multivariate-goodness-of-fit}{%
\subsection{Multivariate Goodness of
Fit}\label{multivariate-goodness-of-fit}}

We can access the individual components of \texttt{results} by name
(\texttt{dir(results)} shows us what is available). Hence
\texttt{results.rsquared} gives us the \(R^2\), and
\texttt{np.sqrt(results.scale)} gives us the RSE.

Variance inflation factors (section 3.3.3) are sometimes useful to
assess the effect of collinearity in the model matrix of a regression
model. We will compute the VIFs in our multiple regression fit, and use
the opportunity to introduce the idea of \emph{list comprehension}.

\hypertarget{list-comprehension}{%
\subsubsection{List Comprehension}\label{list-comprehension}}

Often we encounter a sequence of objects which we would like to
transform for some other task. Below, we compute the VIF for each
feature in our \texttt{X} matrix and produce a data frame whose index
agrees with the columns of \texttt{X}. The notion of list comprehension
can often make such a task easier.

List comprehensions are simple and powerful ways to form lists of
\texttt{Python} objects. The language also supports dictionary and
\emph{generator} comprehension, though these are beyond our scope here.
Let's look at an example. We compute the VIF for each of the variables
in the model matrix \texttt{X}, using the function
\texttt{variance\_inflation\_factor()}.

    \begin{tcolorbox}[breakable, size=fbox, boxrule=1pt, pad at break*=1mm,colback=cellbackground, colframe=cellborder]
\prompt{In}{incolor}{29}{\boxspacing}
\begin{Verbatim}[commandchars=\\\{\}]
\PY{n}{vals} \PY{o}{=} \PY{p}{[}\PY{n}{VIF}\PY{p}{(}\PY{n}{X}\PY{p}{,} \PY{n}{i}\PY{p}{)}
        \PY{k}{for} \PY{n}{i} \PY{o+ow}{in} \PY{n+nb}{range}\PY{p}{(}\PY{l+m+mi}{1}\PY{p}{,} \PY{n}{X}\PY{o}{.}\PY{n}{shape}\PY{p}{[}\PY{l+m+mi}{1}\PY{p}{]}\PY{p}{)}\PY{p}{]}
\PY{n}{vif} \PY{o}{=} \PY{n}{pd}\PY{o}{.}\PY{n}{DataFrame}\PY{p}{(}\PY{p}{\PYZob{}}\PY{l+s+s1}{\PYZsq{}}\PY{l+s+s1}{vif}\PY{l+s+s1}{\PYZsq{}}\PY{p}{:}\PY{n}{vals}\PY{p}{\PYZcb{}}\PY{p}{,}
                   \PY{n}{index}\PY{o}{=}\PY{n}{X}\PY{o}{.}\PY{n}{columns}\PY{p}{[}\PY{l+m+mi}{1}\PY{p}{:}\PY{p}{]}\PY{p}{)}
\PY{n}{vif}
\end{Verbatim}
\end{tcolorbox}

            \begin{tcolorbox}[breakable, size=fbox, boxrule=.5pt, pad at break*=1mm, opacityfill=0]
\prompt{Out}{outcolor}{29}{\boxspacing}
\begin{Verbatim}[commandchars=\\\{\}]
              vif
crim     1.767486
zn       2.298459
indus    3.987181
chas     1.071168
nox      4.369093
rm       1.912532
age      3.088232
dis      3.954037
rad      7.445301
tax      9.002158
ptratio  1.797060
lstat    2.870777
\end{Verbatim}
\end{tcolorbox}
        
    The function \texttt{VIF()} takes two arguments: a dataframe or array,
and a variable column index. In the code above we call \texttt{VIF()} on
the fly for all columns in \texttt{X}.\\
We have excluded column 0 above (the intercept), which is not of
interest. In this case the VIFs are not that exciting.

The object \texttt{vals} above could have been constructed with the
following for loop:

    \begin{tcolorbox}[breakable, size=fbox, boxrule=1pt, pad at break*=1mm,colback=cellbackground, colframe=cellborder]
\prompt{In}{incolor}{30}{\boxspacing}
\begin{Verbatim}[commandchars=\\\{\}]
\PY{n}{vals} \PY{o}{=} \PY{p}{[}\PY{p}{]}
\PY{k}{for} \PY{n}{i} \PY{o+ow}{in} \PY{n+nb}{range}\PY{p}{(}\PY{l+m+mi}{1}\PY{p}{,} \PY{n}{X}\PY{o}{.}\PY{n}{values}\PY{o}{.}\PY{n}{shape}\PY{p}{[}\PY{l+m+mi}{1}\PY{p}{]}\PY{p}{)}\PY{p}{:}
    \PY{n}{vals}\PY{o}{.}\PY{n}{append}\PY{p}{(}\PY{n}{VIF}\PY{p}{(}\PY{n}{X}\PY{o}{.}\PY{n}{values}\PY{p}{,} \PY{n}{i}\PY{p}{)}\PY{p}{)}
\end{Verbatim}
\end{tcolorbox}

    List comprehension allows us to perform such repetitive operations in a
more straightforward way. \#\# Interaction Terms It is easy to include
interaction terms in a linear model using \texttt{ModelSpec()}.
Including a tuple \texttt{("lstat","age")} tells the model matrix
builder to include an interaction term between \texttt{lstat} and
\texttt{age}.

    \begin{tcolorbox}[breakable, size=fbox, boxrule=1pt, pad at break*=1mm,colback=cellbackground, colframe=cellborder]
\prompt{In}{incolor}{31}{\boxspacing}
\begin{Verbatim}[commandchars=\\\{\}]
\PY{n}{X} \PY{o}{=} \PY{n}{MS}\PY{p}{(}\PY{p}{[}\PY{l+s+s1}{\PYZsq{}}\PY{l+s+s1}{lstat}\PY{l+s+s1}{\PYZsq{}}\PY{p}{,}
        \PY{l+s+s1}{\PYZsq{}}\PY{l+s+s1}{age}\PY{l+s+s1}{\PYZsq{}}\PY{p}{,}
        \PY{p}{(}\PY{l+s+s1}{\PYZsq{}}\PY{l+s+s1}{lstat}\PY{l+s+s1}{\PYZsq{}}\PY{p}{,} \PY{l+s+s1}{\PYZsq{}}\PY{l+s+s1}{age}\PY{l+s+s1}{\PYZsq{}}\PY{p}{)}\PY{p}{]}\PY{p}{)}\PY{o}{.}\PY{n}{fit\PYZus{}transform}\PY{p}{(}\PY{n}{Boston}\PY{p}{)}
\PY{n}{model2} \PY{o}{=} \PY{n}{sm}\PY{o}{.}\PY{n}{OLS}\PY{p}{(}\PY{n}{y}\PY{p}{,} \PY{n}{X}\PY{p}{)}
\PY{n}{summarize}\PY{p}{(}\PY{n}{model2}\PY{o}{.}\PY{n}{fit}\PY{p}{(}\PY{p}{)}\PY{p}{)}
\end{Verbatim}
\end{tcolorbox}

            \begin{tcolorbox}[breakable, size=fbox, boxrule=.5pt, pad at break*=1mm, opacityfill=0]
\prompt{Out}{outcolor}{31}{\boxspacing}
\begin{Verbatim}[commandchars=\\\{\}]
              coef  std err       t  P>|t|
intercept  36.0885    1.470  24.553  0.000
lstat      -1.3921    0.167  -8.313  0.000
age        -0.0007    0.020  -0.036  0.971
lstat:age   0.0042    0.002   2.244  0.025
\end{Verbatim}
\end{tcolorbox}
        
    \hypertarget{non-linear-transformations-of-the-predictors}{%
\subsection{Non-linear Transformations of the
Predictors}\label{non-linear-transformations-of-the-predictors}}

The model matrix builder can include terms beyond just column names and
interactions. For instance, the \texttt{poly()} function supplied in
\texttt{ISLP} specifies that columns representing polynomial functions
of its first argument are added to the model matrix.

    \begin{tcolorbox}[breakable, size=fbox, boxrule=1pt, pad at break*=1mm,colback=cellbackground, colframe=cellborder]
\prompt{In}{incolor}{32}{\boxspacing}
\begin{Verbatim}[commandchars=\\\{\}]
\PY{n}{X} \PY{o}{=} \PY{n}{MS}\PY{p}{(}\PY{p}{[}\PY{n}{poly}\PY{p}{(}\PY{l+s+s1}{\PYZsq{}}\PY{l+s+s1}{lstat}\PY{l+s+s1}{\PYZsq{}}\PY{p}{,} \PY{n}{degree}\PY{o}{=}\PY{l+m+mi}{2}\PY{p}{)}\PY{p}{,} \PY{l+s+s1}{\PYZsq{}}\PY{l+s+s1}{age}\PY{l+s+s1}{\PYZsq{}}\PY{p}{]}\PY{p}{)}\PY{o}{.}\PY{n}{fit\PYZus{}transform}\PY{p}{(}\PY{n}{Boston}\PY{p}{)}
\PY{n}{model3} \PY{o}{=} \PY{n}{sm}\PY{o}{.}\PY{n}{OLS}\PY{p}{(}\PY{n}{y}\PY{p}{,} \PY{n}{X}\PY{p}{)}
\PY{n}{results3} \PY{o}{=} \PY{n}{model3}\PY{o}{.}\PY{n}{fit}\PY{p}{(}\PY{p}{)}
\PY{n}{summarize}\PY{p}{(}\PY{n}{results3}\PY{p}{)}
\end{Verbatim}
\end{tcolorbox}

            \begin{tcolorbox}[breakable, size=fbox, boxrule=.5pt, pad at break*=1mm, opacityfill=0]
\prompt{Out}{outcolor}{32}{\boxspacing}
\begin{Verbatim}[commandchars=\\\{\}]
                              coef  std err       t  P>|t|
intercept                  17.7151    0.781  22.681    0.0
poly(lstat, degree=2)[0] -179.2279    6.733 -26.620    0.0
poly(lstat, degree=2)[1]   72.9908    5.482  13.315    0.0
age                         0.0703    0.011   6.471    0.0
\end{Verbatim}
\end{tcolorbox}
        
    The effectively zero \emph{p}-value associated with the quadratic term
(i.e.~the third row above) suggests that it leads to an improved model.

By default, \texttt{poly()} creates a basis matrix for inclusion in the
model matrix whose columns are \emph{orthogonal polynomials}, which are
designed for stable least squares computations. \{Actually,
\texttt{poly()} is a wrapper for the workhorse and standalone function
\texttt{Poly()} that does the work in building the model matrix.\}
Alternatively, had we included an argument \texttt{raw=True} in the
above call to \texttt{poly()}, the basis matrix would consist simply of
\texttt{lstat} and \texttt{lstat**2}. Since either of these bases
represent quadratic polynomials, the fitted values would not change in
this case, just the polynomial coefficients. Also by default, the
columns created by \texttt{poly()} do not include an intercept column as
that is automatically added by \texttt{MS()}.

We use the \texttt{anova\_lm()} function to further quantify the extent
to which the quadratic fit is superior to the linear fit.

    \begin{tcolorbox}[breakable, size=fbox, boxrule=1pt, pad at break*=1mm,colback=cellbackground, colframe=cellborder]
\prompt{In}{incolor}{33}{\boxspacing}
\begin{Verbatim}[commandchars=\\\{\}]
\PY{n}{anova\PYZus{}lm}\PY{p}{(}\PY{n}{results1}\PY{p}{,} \PY{n}{results3}\PY{p}{)}
\end{Verbatim}
\end{tcolorbox}

            \begin{tcolorbox}[breakable, size=fbox, boxrule=.5pt, pad at break*=1mm, opacityfill=0]
\prompt{Out}{outcolor}{33}{\boxspacing}
\begin{Verbatim}[commandchars=\\\{\}]
   df\_resid           ssr  df\_diff      ss\_diff           F        Pr(>F)
0     503.0  19168.128609      0.0          NaN         NaN           NaN
1     502.0  14165.613251      1.0  5002.515357  177.278785  7.468491e-35
\end{Verbatim}
\end{tcolorbox}
        
    Here \texttt{results1} represents the linear submodel containing
predictors \texttt{lstat} and \texttt{age}, while \texttt{results3}
corresponds to the larger model above with a quadratic term in
\texttt{lstat}. The \texttt{anova\_lm()} function performs a hypothesis
test comparing the two models. The null hypothesis is that the quadratic
term in the bigger model is not needed, and the alternative hypothesis
is that the bigger model is superior. Here the \emph{F}-statistic is
177.28 and the associated \emph{p}-value is zero. In this case the
\emph{F}-statistic is the square of the \emph{t}-statistic for the
quadratic term in the linear model summary for \texttt{results3} --- a
consequence of the fact that these nested models differ by one degree of
freedom. This provides very clear evidence that the quadratic polynomial
in \texttt{lstat} improves the linear model. This is not surprising,
since earlier we saw evidence for non-linearity in the relationship
between \texttt{medv} and \texttt{lstat}.

The function \texttt{anova\_lm()} can take more than two nested models
as input, in which case it compares every successive pair of models.
That also explains why there are \texttt{NaN}s in the first row above,
since there is no previous model with which to compare the first.

    \begin{tcolorbox}[breakable, size=fbox, boxrule=1pt, pad at break*=1mm,colback=cellbackground, colframe=cellborder]
\prompt{In}{incolor}{34}{\boxspacing}
\begin{Verbatim}[commandchars=\\\{\}]
\PY{n}{ax} \PY{o}{=} \PY{n}{subplots}\PY{p}{(}\PY{n}{figsize}\PY{o}{=}\PY{p}{(}\PY{l+m+mi}{8}\PY{p}{,}\PY{l+m+mi}{8}\PY{p}{)}\PY{p}{)}\PY{p}{[}\PY{l+m+mi}{1}\PY{p}{]}
\PY{n}{ax}\PY{o}{.}\PY{n}{scatter}\PY{p}{(}\PY{n}{results3}\PY{o}{.}\PY{n}{fittedvalues}\PY{p}{,} \PY{n}{results3}\PY{o}{.}\PY{n}{resid}\PY{p}{)}
\PY{n}{ax}\PY{o}{.}\PY{n}{set\PYZus{}xlabel}\PY{p}{(}\PY{l+s+s1}{\PYZsq{}}\PY{l+s+s1}{Fitted value}\PY{l+s+s1}{\PYZsq{}}\PY{p}{)}
\PY{n}{ax}\PY{o}{.}\PY{n}{set\PYZus{}ylabel}\PY{p}{(}\PY{l+s+s1}{\PYZsq{}}\PY{l+s+s1}{Residual}\PY{l+s+s1}{\PYZsq{}}\PY{p}{)}
\PY{n}{ax}\PY{o}{.}\PY{n}{axhline}\PY{p}{(}\PY{l+m+mi}{0}\PY{p}{,} \PY{n}{c}\PY{o}{=}\PY{l+s+s1}{\PYZsq{}}\PY{l+s+s1}{k}\PY{l+s+s1}{\PYZsq{}}\PY{p}{,} \PY{n}{ls}\PY{o}{=}\PY{l+s+s1}{\PYZsq{}}\PY{l+s+s1}{\PYZhy{}\PYZhy{}}\PY{l+s+s1}{\PYZsq{}}\PY{p}{)}\PY{p}{;}
\end{Verbatim}
\end{tcolorbox}

    \begin{center}
    \adjustimage{max size={0.9\linewidth}{0.9\paperheight}}{artifacts/latex/Ch03-linreg-lab_files/artifacts/latex/Ch03-linreg-lab_73_0.png}
    \end{center}
    { \hspace*{\fill} \\}
    
    We see that when the quadratic term is included in the model, there is
little discernible pattern in the residuals. In order to create a cubic
or higher-degree polynomial fit, we can simply change the degree
argument to \texttt{poly()}.

    \hypertarget{qualitative-predictors}{%
\subsection{Qualitative Predictors}\label{qualitative-predictors}}

Here we use the \texttt{Carseats} data, which is included in the
\texttt{ISLP} package. We will attempt to predict \texttt{Sales} (child
car seat sales) in 400 locations based on a number of predictors.

    \begin{tcolorbox}[breakable, size=fbox, boxrule=1pt, pad at break*=1mm,colback=cellbackground, colframe=cellborder]
\prompt{In}{incolor}{35}{\boxspacing}
\begin{Verbatim}[commandchars=\\\{\}]
\PY{n}{Carseats} \PY{o}{=} \PY{n}{load\PYZus{}data}\PY{p}{(}\PY{l+s+s1}{\PYZsq{}}\PY{l+s+s1}{Carseats}\PY{l+s+s1}{\PYZsq{}}\PY{p}{)}
\PY{n}{Carseats}\PY{o}{.}\PY{n}{columns}
\end{Verbatim}
\end{tcolorbox}

            \begin{tcolorbox}[breakable, size=fbox, boxrule=.5pt, pad at break*=1mm, opacityfill=0]
\prompt{Out}{outcolor}{35}{\boxspacing}
\begin{Verbatim}[commandchars=\\\{\}]
Index(['Sales', 'CompPrice', 'Income', 'Advertising', 'Population', 'Price',
       'ShelveLoc', 'Age', 'Education', 'Urban', 'US'],
      dtype='object')
\end{Verbatim}
\end{tcolorbox}
        
    The \texttt{Carseats}\\
data includes qualitative predictors such as \texttt{ShelveLoc}, an
indicator of the quality of the shelving location --- that is, the space
within a store in which the car seat is displayed. The predictor
\texttt{ShelveLoc} takes on three possible values, \texttt{Bad},
\texttt{Medium}, and \texttt{Good}. Given a qualitative variable such as
\texttt{ShelveLoc}, \texttt{ModelSpec()} generates dummy variables
automatically. These variables are often referred to as a \emph{one-hot
encoding} of the categorical feature. Their columns sum to one, so to
avoid collinearity with an intercept, the first column is dropped. Below
we see the column \texttt{ShelveLoc{[}Bad{]}} has been dropped, since
\texttt{Bad} is the first level of \texttt{ShelveLoc}. Below we fit a
multiple regression model that includes some interaction terms.

    \begin{tcolorbox}[breakable, size=fbox, boxrule=1pt, pad at break*=1mm,colback=cellbackground, colframe=cellborder]
\prompt{In}{incolor}{36}{\boxspacing}
\begin{Verbatim}[commandchars=\\\{\}]
\PY{n}{allvars} \PY{o}{=} \PY{n+nb}{list}\PY{p}{(}\PY{n}{Carseats}\PY{o}{.}\PY{n}{columns}\PY{o}{.}\PY{n}{drop}\PY{p}{(}\PY{l+s+s1}{\PYZsq{}}\PY{l+s+s1}{Sales}\PY{l+s+s1}{\PYZsq{}}\PY{p}{)}\PY{p}{)}
\PY{n}{y} \PY{o}{=} \PY{n}{Carseats}\PY{p}{[}\PY{l+s+s1}{\PYZsq{}}\PY{l+s+s1}{Sales}\PY{l+s+s1}{\PYZsq{}}\PY{p}{]}
\PY{n}{final} \PY{o}{=} \PY{n}{allvars} \PY{o}{+} \PY{p}{[}\PY{p}{(}\PY{l+s+s1}{\PYZsq{}}\PY{l+s+s1}{Income}\PY{l+s+s1}{\PYZsq{}}\PY{p}{,} \PY{l+s+s1}{\PYZsq{}}\PY{l+s+s1}{Advertising}\PY{l+s+s1}{\PYZsq{}}\PY{p}{)}\PY{p}{,}
                   \PY{p}{(}\PY{l+s+s1}{\PYZsq{}}\PY{l+s+s1}{Price}\PY{l+s+s1}{\PYZsq{}}\PY{p}{,} \PY{l+s+s1}{\PYZsq{}}\PY{l+s+s1}{Age}\PY{l+s+s1}{\PYZsq{}}\PY{p}{)}\PY{p}{]}
\PY{n}{X} \PY{o}{=} \PY{n}{MS}\PY{p}{(}\PY{n}{final}\PY{p}{)}\PY{o}{.}\PY{n}{fit\PYZus{}transform}\PY{p}{(}\PY{n}{Carseats}\PY{p}{)}
\PY{n}{model} \PY{o}{=} \PY{n}{sm}\PY{o}{.}\PY{n}{OLS}\PY{p}{(}\PY{n}{y}\PY{p}{,} \PY{n}{X}\PY{p}{)}
\PY{n}{summarize}\PY{p}{(}\PY{n}{model}\PY{o}{.}\PY{n}{fit}\PY{p}{(}\PY{p}{)}\PY{p}{)}
\end{Verbatim}
\end{tcolorbox}

            \begin{tcolorbox}[breakable, size=fbox, boxrule=.5pt, pad at break*=1mm, opacityfill=0]
\prompt{Out}{outcolor}{36}{\boxspacing}
\begin{Verbatim}[commandchars=\\\{\}]
                      coef  std err       t  P>|t|
intercept           6.5756    1.009   6.519  0.000
CompPrice           0.0929    0.004  22.567  0.000
Income              0.0109    0.003   4.183  0.000
Advertising         0.0702    0.023   3.107  0.002
Population          0.0002    0.000   0.433  0.665
Price              -0.1008    0.007 -13.549  0.000
ShelveLoc[Good]     4.8487    0.153  31.724  0.000
ShelveLoc[Medium]   1.9533    0.126  15.531  0.000
Age                -0.0579    0.016  -3.633  0.000
Education          -0.0209    0.020  -1.063  0.288
Urban[Yes]          0.1402    0.112   1.247  0.213
US[Yes]            -0.1576    0.149  -1.058  0.291
Income:Advertising  0.0008    0.000   2.698  0.007
Price:Age           0.0001    0.000   0.801  0.424
\end{Verbatim}
\end{tcolorbox}
        
    In the first line above, we made \texttt{allvars} a list, so that we
could add the interaction terms two lines down. Our model-matrix builder
has created a \texttt{ShelveLoc{[}Good{]}} dummy variable that takes on
a value of 1 if the shelving location is good, and 0 otherwise. It has
also created a \texttt{ShelveLoc{[}Medium{]}} dummy variable that equals
1 if the shelving location is medium, and 0 otherwise. A bad shelving
location corresponds to a zero for each of the two dummy variables. The
fact that the coefficient for \texttt{ShelveLoc{[}Good{]}} in the
regression output is positive indicates that a good shelving location is
associated with high sales (relative to a bad location). And
\texttt{ShelveLoc{[}Medium{]}} has a smaller positive coefficient,
indicating that a medium shelving location leads to higher sales than a
bad shelving location, but lower sales than a good shelving location.


    % Add a bibliography block to the postdoc
    
    
    
\end{document}
